% !TEX program = xelatex
\documentclass[11pt,a4paper]{article}

\usepackage[UTF8]{ctex}
\usepackage[english]{babel}
\usepackage[margin=2.5cm]{geometry}
\usepackage{setspace}
\onehalfspacing

\usepackage{amsmath,amssymb,amsthm}
\usepackage{mathtools}
\usepackage{bm}

\newtheorem{theorem}{定理}[section]
\newtheorem{lemma}[theorem]{引理}
\newtheorem{proposition}[theorem]{命题}
\newtheorem{corollary}[theorem]{推论}
\newtheorem{definition}[theorem]{定义}
\newtheorem{remark}[theorem]{注记}
\theoremstyle{definition}
\newtheorem{condition}[theorem]{条件}

\usepackage{graphicx}
\usepackage{float}
\usepackage{booktabs}
\usepackage{multirow}
\usepackage{array}
\usepackage{caption}
\usepackage{subcaption}

\usepackage[colorlinks=true,linkcolor=blue,citecolor=blue,urlcolor=blue]{hyperref}

\usepackage{csquotes}
\usepackage[style=authoryear-comp,backend=biber]{biblatex}
\addbibresource{references.bib}

\usepackage{enumitem}
\usepackage{tikz}
\usetikzlibrary{arrows.meta,positioning,shapes,fit,backgrounds}
\usepackage{mdframed}
\usepackage{xcolor}

\definecolor{lightgray}{RGB}{245,245,245}
\definecolor{lightblue}{RGB}{230,240,250}
\definecolor{lightred}{RGB}{255,240,240}

\newcommand{\verifspace}{\mathcal{C}_{\mathcal{V}}}
\newcommand{\taskcspace}{\mathcal{C}}
\newcommand{\effcap}{\mathcal{E}}
\newcommand{\contextcap}{\mathcal{W}}

\title{{\bfseries 收敛与发散：\\LLM Agent的能力边界、结构性缺失与人机协作的必然性}}

\author{}

\date{\today}

\begin{document}
\maketitle

\begin{abstract}
\noindent
当前基于大语言模型的Agent系统在工程实践中反复出现一组稳定的失败模式——多步规划中的语义漂移、全自动闭环的不可靠性、跨会话记忆的结构性退化——这些模式随模型代际提升而减轻，却未被消除。主流研究倾向于将这些失败归因于噪声积累、上下文过载或规划-执行脱节等工程层面的不足，并诉诸更大规模模型、更长上下文窗口或更精细的提示词工程来修补。本文提出一种根本不同的立场：上述失败中相当一部分不是待解决的工程问题，而是从当前Transformer架构的数学结构中推导出的结构必然。

本文的论证路径如下。首先，基于残差流张量积 $\mathbb{R}^n \otimes \mathbb{R}^d$ 的结构分解，建立语义空间（$\mathbb{R}^d$因子，范畴闭合）与路由空间（$\mathbb{R}^n$因子，跨层范畴不闭合）的精确定位，由此导出不对称闭合——$\mathbb{R}^d$的W矩阵可复合、$\mathbb{R}^n$的A矩阵不可复合、且前者的复合成功是后者复合失败的因果根源——作为全文最底层的数学基础。在不对称闭合之上，建立训练收敛与测试收敛的双模分类框架，论证记忆中不存在长短期之间的连续过渡结构，并确认参数空间有效维度的截断投影限制。进而引入边缘混沌假说——不复合导致的路径依赖加非线性折叠加局部扩张，使得训练良好的Transformer处于最大Lyapunov指数趋近于零的临界动力学状态——为语义聚类的涌现提供统一的动力学解释。在此基础上，将Agent行为按评价模式二分为收敛任务（具备可自动执行的验证标准）与发散任务（验证标准不存在或不可形式化）。对收敛任务，以三段论形式论证上下文驱动分解的结构不可靠性，并导出层级收敛结构中根需求必须由外部锚定的闭合限制。对发散任务，将"探索"还原为条件概率采样机制，建立偏好驱动迭代理论。最后，从Agent相对于人的三项结构性缺失出发，系统推导人在人机协作中的不可替代角色。

本文的主要贡献包括：(1)将残差流 $\mathbb{R}^n \otimes \mathbb{R}^d$ 分解为语义空间与路由空间两个正交且因果耦合的因子，精确定位了Transformer中范畴结构的闭合与不闭合区域，将传统上被视为"弱闭合"缺陷的不对称性重新理解为语义动力学的前提条件；(2)提出边缘混沌假说，将注意力不复合的动力学后果统一刻画，并给出可检验的预测；(3)提出训练收敛与测试收敛的双模分类；(4)以三段论形式证明上下文驱动分解决策的不可靠性；(5)建立发散任务的偏好驱动迭代理论；(6)从Agent的结构性缺失中推导人在人机协作中的构成性角色。

全文采用六层论证强度分层——严格层、定义层、形式论证层、概念构造层、类比层、经验推论层——各节所属层级在正文中明确标注。
\end{abstract}

\newpage
\tableofcontents
\newpage

% ===================================================================
\section{引言：从使用原则到理论动机}
\label{sec:intro}
% ===================================================================

\subsection{经验归纳：十七项使用原则}

以下原则来自本文作者在长期工程实践中对LLM Agent行为的系统性观察与归纳。这些原则在实践中反复验证有效，但其有效性背后的统一机制尚未被现有理论充分解释。本节以命题形式陈述这些原则，作为全文理论构建的经验起点。为便于后续各节的论证引用，每条原则赋予一个编号（P1--P17）。

\medskip
\textbf{第一组：输入原则（P1--P3）}

\textbf{P1（精确优先于简短）}：提示词的信息角色明确性优先于文本长度最小化。有限维精确输入降低信息的结构复杂度，提高输入落入收敛覆盖区域的概率。

\textbf{P2（词汇的分布场效应）}：词汇在模型中的效力不源于其字面语义定义，而源于该词汇在训练数据中典型出现的语用场景。要求模型保持批判性，使用"批判""实事求是"等常出现在严格要求语境中的词汇，其效果优于直接使用抽象标签"客观"。这一现象是分布语义学原理——在范畴论中即米田引理（Yoneda Lemma）：在FinStoch范畴中，一个对象的语义身份由所有指向它的态射的集合所唯一确定——的直接体现。

\textbf{P3（目标约束而非路径约束）}：向模型交代目标、边界条件和验证标准，而不规定具体实现步骤。每增加一个路径层面的指令，即增加一维潜在的语义歧义，歧义维度在条件概率空间中以指数方式放大偏离收敛覆盖区域的概率。

\medskip
\textbf{第二组：上下文管理原则（P4--P5）}

\textbf{P4（上下文信息的审慎选择）}：注入上下文窗口的每条信息均需审慎评估其信息论价值。模型的有效上下文容量显著小于名义窗口长度；注意力机制在结构上对所有窗口内信息施加等权覆盖，冗余信息与关键信息在注意力计算中无优先级区分地竞争有限的表示资源。

\textbf{P5（上下文的分级滚动管理）}：上下文内容按信息优先级进行分级组织，高优先级信息（任务目标、约束条件、已验证的中间结果）在窗口内常驻保留，低优先级信息随窗口推进按优先级梯度滚动淘汰。此策略是在有限有效容量约束下的信息论最优分配。

\medskip
\textbf{第三组：收敛任务原则（P6--P10）}

\textbf{P6（能力边界的节点划分）}：在模型文字输入的有效边界处对任务执行节点切分。超出有效边界的信息不受训练收敛约束的覆盖，切分回边界内方能恢复收敛约束对映射行为的锻造作用，使测试收敛获得运作空间。

\textbf{P7（测试收敛驱动）}：Agent在运行时的行为不由梯度下降驱动，而由"生成$\to$验证$\to$分析$\to$修正"的迭代循环驱动。参数在训练完成后静态冻结，训练收敛通道关闭；Agent行为仅能由可验证目标驱动的测试收敛来导向。

\textbf{P8（外部约束锚定）}：任务分解中的关键分叉决策须由独立于模型的硬约束——人工定义或程序逻辑——进行仲裁，不得由模型在上下文中自行判断。节点错误沿层级向下乘性传导，下层子任务在偏差基底上执行且缺乏独立纠正偏差所需的信息。

\textbf{P9（生成与检验的分离）}：生成输出与检验输出不得由同一模型实例在同一上下文窗口中完成。检验过程本身是上下文中的推理，同样未经历收敛约束——自写自验构成循环盲区，检验的有效性无法独立于被检验对象而确立。

\textbf{P10（逐环节验证）}：任务层级结构中每个子环节完成即进行独立验证，不得累积至末端统一检验。前一环节的未验证输出若作为后续环节的输入参与计算，错误将沿层级乘性放大——早期验证即截断错误传播链。

\medskip
\textbf{第四组：发散任务原则（P11--P12）}

\textbf{P11（发散与收敛的空间隔离）}：发散探索与收敛执行须使用相互独立的上下文空间，二者之间仅通过可验证目标结晶后的显式过渡协议连接。发散阶段累积的偏好标注——包含被淘汰的候选方向和已做出的偏好判断——与收敛阶段的验证状态若共享上下文，将产生双向污染：偏好信息干扰收敛条件，验证逻辑压制探索空间。

\textbf{P12（偏好驱动的对话式协作）}：发散任务中，人持续提供比较性偏好信号，模型在偏好信号对条件分布的连续重塑下生成候选方案。发散任务不具备可自动执行的验证标准，人的偏好判断是采样器在语义空间中唯一的方向梯度来源。

\medskip
\textbf{第五组：结构边界原则（P13--P17）}

\textbf{P13（扁平化管理）}：子Agent的管理结构采用单层级扁平模式而非多层嵌套。每增加一个嵌套层级，该层级的驱动指令（由上一层Agent在上下文中生成）的不可靠性沿层级乘性放大，高层收敛成本累积将超出上下文窗口的承载极限。

\textbf{P14（操作者上限）}：操作者对任务结构的理解深度和判断质量直接构成Agent在该任务上实际可达到的性能上限。根需求锚定、关键分叉仲裁、偏好信号提供——每一项Agent结构上不可自举的环节均由操作者填补，操作者的认知边界即系统在该任务上的边界。

\textbf{P15（验证覆盖即能力边界）}：Agent在给定任务上可达到的可靠性上界，不超过该任务可被自动检验的约束维度范围。测试收敛的有效性等价于验证信号对任务需求语义空间的覆盖度——未被验证信号覆盖的约束维度上，Agent可以系统性偏离而不触发任何纠正信号。

\textbf{P16（根需求的外在性）}：任务链中最高层的需求目标必须由系统外部（操作者）提供并锚定。验证标准沿层级递归上溯，终点为"当前产出是否符合用户的真实意图"——该判断的语义内容位于模型训练分布之外，不在收敛覆盖范围之内。全自动闭环在根节点处原则上不可闭合。

\textbf{P17（幻觉边界不可消除）}：模型在收敛覆盖范围之外必然产生不受收敛约束的输出。训练改进和微调可以扩展覆盖范围——使受约束的区域面积增大——但边界本身来自参数空间有限维截断的数学结构，无法被消除。

\subsection{核心问题}

上述十七条原则共同指向一个结构根源：\textbf{在LLM Agent的信息处理链路中，存在一条根本性的二分线——某些信息经历了训练收敛的梯度反馈锻造，某些信息从未经历。}模型参数经历了这一锻造过程；上下文窗口中的信息则直接参与前向计算而未受梯度筛选。这一区分来自Transformer架构的底层数学结构——具体而言，来自残差流 $\mathbb{R}^n \otimes \mathbb{R}^d$ 中语义空间（$\mathbb{R}^d$）与路由空间（$\mathbb{R}^n$）两个因子的范畴状态不对称：一个闭合，一个不闭合。这一结构性区分是否可消除？若不可消除，它如何系统性地决定了Agent在不同任务条件下的能力边界？本文以这两个问题为出发点，逐层展开论证。

\subsection{方法论立场}

本文的工作属于\textbf{概念闭合、证明分层的框架构建传统}。该传统在理论物理学和系统理论中具有长期先例——从少数第一性假设出发，以分层论证的方式推导覆盖广域的结论体系，在每一论证层级诚实标注论证的完成程度与未闭合的缺口。Rosen (1991)的"生命系统范畴论"和Maturana与Varela (1980)的"自创生理论"均体现了这一方法论：概念体系完整且自洽，数学形式化按可及程度分层推进，未完成的部分被标注为开放问题而非论证缺陷。近年AI研究中的先例包括Chollet (2019)对智能定义的形式化框架和Legg与Hutter (2007)的通用智能度量——二者均采用"核心部分形式化、外围部分概念构造、边界处诚实标注不可计算性或未闭合性"的论证结构。

本文的目标不是提供一套已完成的形式化定理体系，而是完成以下三步论证：(1)从当前Transformer架构的少数数学结构约束出发——残差流的 $\mathbb{R}^n \otimes \mathbb{R}^d$ 张量积结构、语义空间与路由空间的范畴状态不对称、参数空间的有限维截断、上下文信息不经训练收敛、参数化记忆与序列化记忆的数学结构异质性——可以推导出LLM Agent在收敛任务和发散任务中的系统性行为边界；(2)这些行为边界构成了工程实践中最有效的使用原则——包括第\ref{sec:intro}节中归纳的十七条原则——背后统一的结构解释；(3)未被当前数学结构闭合的部分——包括范畴论强闭合的形式化、发散任务中偏好信号的量化理论、混合任务的复杂性分析、边缘混沌假说的经验验证——被诚实标注为开放问题，构成后续研究的路线图。其中，第(1)部分中已完成严格形式论证的环节（多级需求收敛结构的推导）构成框架的逻辑内核；以概念构造和功能类比方式建立的环节构成框架的解释性外延；诚实标注的未闭合项构成框架的自我修正机制。


% ===================================================================
\section{第一性原理理论结构}
\label{sec:foundations}

回答第\ref{sec:intro}节的核心追问需要从Transformer最底层的数学对象——残差流——出发。本节以残差流 $\mathbb{R}^n \otimes \mathbb{R}^d$ 的结构分解为起点，区分语义空间与路由空间两个正交且因果耦合的因子，建立不对称闭合的精确图景，并在此基础上依次展开收敛二分、截断投影、记忆二分、边缘混沌、以及使用原则的统一推导。论证强度：\textbf{严格层}覆盖2.1--2.4；\textbf{形式论证层}覆盖2.5--2.6；\textbf{概念构造层}覆盖2.7；\textbf{类比层}覆盖2.8。

% -------------------------------------------------------------------
\subsection{残差流的数学结构：$\mathbb{R}^n \otimes \mathbb{R}^d$ 分解}
\label{sec:residual_decomp}

Transformer每一层的状态是一个矩阵：

\begin{equation}
X^{(l)} \in \mathbb{R}^n \otimes \mathbb{R}^d
\end{equation}

其中 $n$ 为token数量，$d$ 为特征维度。$X^{(l)}$ 的第 $i$ 行 $x_i^{(l)} \in \mathbb{R}^d$ 是token $i$ 在当前层的语义向量。\textbf{同一个矩阵 $X^{(l)}$ 携带两种正交且因果耦合的信息}——$\mathbb{R}^d$ 因子承载语义（每行向量的含义），$\mathbb{R}^n$ 因子承载路由（行与行之间的注意力权重）。张量积结构意味着两个因子的运算是独立的——注意力A在行方向混合token（$\mathbb{R}^n$上的运算），权重矩阵W在列方向变换特征（$\mathbb{R}^d$上的运算），先做哪个后做哪个互不干涉。但正交不等于无关——A的条目依赖 $x \in \mathbb{R}^d$ 的值，W改变 $x$，所以W的输出影响下一层A的输入。两者构成\textbf{反馈环}：W改变语义坐标 $\to$ A的路由变化 $\to$ 新的语义坐标 $\to$ 下一层W作用于新坐标 $\to$ A的路由再次变化。以下三个子节分别考察两个因子的范畴状态及它们的因果耦合。

\subsubsection{语义空间 $\mathbb{R}^d$：范畴闭合}
\label{sec:semantic_space}

$\mathbb{R}^d$ 上的操作由一组固定的线性算子——权重矩阵 $W_Q, W_K, W_V, W_O$ 和MLP矩阵——完成。它们不随输入变化，是训练后固定的参数。

\textbf{$\mathbb{R}^d$ 构成范畴（Mat / Vect）：}
\begin{itemize}[leftmargin=*,itemsep=0.2em]
    \item 对象：$\mathbb{R}^d$ 及其子空间
    \item 态射：W矩阵——从 $\mathbb{R}^d$ 到 $\mathbb{R}^d$ 的线性映射
    \item 复合：$W_2 \cdot W_1$ 始终定义良好（矩阵乘法），结合律由矩阵乘法保证
    \item 恒等：单位矩阵 $I_d$
\end{itemize}

\textbf{米田引理在 $\mathbb{R}^d$ 上成立：}语义空间 $\mathbb{R}^d$ 的身份完全由所有指向它的态射 $\hom(-, \mathbb{R}^d)$ 确定——即所有W矩阵及其复合的全体。$\mathbb{R}^d$ 有自足的、不依赖token地址那一端的代数结构。

\textbf{单层注意力是合法的FinStoch态射（数学事实，无近似）：}单头注意力 $A = \text{softmax}(QK^T / \sqrt{d_k})$ 严格产出随机矩阵——非负，行和为1。多头注意力：多个随机矩阵经concat + $W_O$ 做凸组合，随机矩阵集合是凸集，结果仍为随机矩阵。因此单层注意力A是FinStoch中的合法态射——这一步\textbf{没有近似，没有比喻}。米田嵌入的操作化表述：一个token的语义身份由所有其他token到它的条件概率关系所定义——这是分布语义学（Firth 1957："you shall know a word by the company it keeps"）的范畴论精确表述。

训练将全部语义知识压进了 $\mathbb{R}^d$ 的几何中——向量的方向、距离、W矩阵的变换。W矩阵的范畴复合不是"模拟"语义变换——它\textbf{就是}语义变换。Anthropic的虚拟权重和路径展开，做的就是把W的复合作为语义操作来读。

\textbf{论证强度：严格层。}

\subsubsection{路由空间 $\mathbb{R}^n$：跨层范畴不闭合}
\label{sec:routing_space}

注意力矩阵 $A \in [0,1]^{n \times n}$，条目 $a_{ij} = \text{softmax}(Q(x_i) \cdot K(x_j)^T / \sqrt{d})$ 表示"token $i$ 应该从token $j$ 取多少信息"。这是token间的路由决策。

\textbf{单层内：}对固定的表示 $X^{(l)}$，$A^{(l)}$ 是FinStoch态射（行随机矩阵）。成立。

\textbf{跨层：}$A^{(l+1)} \circ A^{(l)}$ \textbf{无定义。}这是全框架最关键的结构事实之一。原因：

\begin{equation}
a_{ij}^{(l)} = f(x_i^{(l)}, x_j^{(l)}; W_Q, W_K) \quad \text{——依赖 $\mathbb{R}^d$ 当前状态}
\end{equation}
\begin{equation}
x^{(l+1)} = x^{(l)} + A^{(l)}V^{(l)}W_O + \text{MLP}(\cdots) \quad \text{——$\mathbb{R}^d$ 状态被改变}
\end{equation}

由此：$A^{(l+1)}$ 的输入在 $\mathbb{R}^d$ 端已被 $A^{(l)}$ 的输出改变 $\to$ $A^{(l+1)}$ 的域对象不是 $A^{(l)}$ 的余域对象 $\to$ Chapman-Kolmogorov方程不成立 $\to$ 复合 $A^{(l+1)} \cdot A^{(l)}$ 不是适当的数学描述。

$\mathbb{R}^n$ 不构成跨层范畴的结构原因：$a_{ij}$ 不由地址 $\{i, j\}$ 决定——它由地址里装的\textbf{语义内容}决定；语义内容每层被W改变 $\to$ $a_{ij}$ 逐层变化；$a_{jk}$ 和 $a_{ij}$ 之间没有运算能得到 $a_{ik}$——不传递；没有恒等态射——自注意力 $a_{ii}$ 做计算，不是"什么都不做"。

\textbf{米田引理不适用：}token位置不构成范畴对象——米田上场的前提不满足。Token $i$ 的身份不能由"指向它的注意力"确定，因为注意力本身依赖 $\mathbb{R}^d$ 端当前状态。Token是空槽位——索引 $i$ 跨层不变（"哪里"），但里面装的向量 $x_i^{(l)}$ 每层被重写（"什么"）。Token没有自足的语义身份——它的身份完全寄生在语义空间 $\mathbb{R}^d$ 上。

\textbf{反事实推理——如果 $\mathbb{R}^n$ 也构成范畴会怎样（关键检验）：}
\begin{quote}
假设米田成立，$a_{ij}$ 由地址 $\{i, j\}$ 自足决定，不依赖 $\mathbb{R}^d$。则A跨层固定 $\to$ $A_{\text{eff}} = A^L$ $\to$ $\{1,\ldots,n\}$ 上的L步Markov链，稠密正矩阵 $\to$ 收敛到唯一平稳分布 $\pi$ $\to$ $A_{\text{eff}}$ 所有行 $= \pi$（所有token注意力分布相同）$\to$ 所有token收到相同的加权上下文 $\to$ 同方向更新 $\to$ 表示同质化 $\to$ 跨token动态路由消失 $\to$ 语义结构丰富性严重退化 $\to$ 深度冗余。
\end{quote}

\textbf{复合闭合 = 跨token动态路由消失 = 语义死亡。}这不是设计疏忽——$\mathbb{R}^n$ 端的不复合是语义聚类能够在 $\mathbb{R}^d$ 端形成的数学边界条件。

\textbf{论证强度：}严格层（跨层复合失败是数学事实）+ 概念构造层（"不复合 = 功能性"的论证）。

\subsubsection{因果耦合与不对称闭合}
\label{sec:causal_coupling}

两个因子的范畴状态不是独立的——它们因果耦合：

\begin{quote}
W矩阵构成范畴 $\to$ W复合改变语义空间坐标 $\to$ 每一层面对的 $\mathbb{R}^d$ 状态不同 \\
$\to$ A矩阵依赖 $\mathbb{R}^d$ 当前状态（$Q \cdot K^T$ 使用当前坐标）\\
$\to$ 语义空间变了 $\to$ A矩阵必然逐层变化 $\to$ A的跨层复合断裂
\end{quote}

\textbf{W的复合成功是A的复合失败的原因。}这不是两个独立事实，是因果关系。只要语义空间在演化（W复合），路由空间就不能累积（A不复合）。

\textbf{不对称闭合：}
\begin{itemize}[leftmargin=*,itemsep=0.2em]
    \item $\mathbb{R}^d$ 端：范畴闭合，复合成立，米田成立 $\to$ 语义有自足结构
    \item $\mathbb{R}^n$ 端：范畴不闭合，复合失败，米田不适用 $\to$ 路由寄生在语义上
    \item 因果方向：$\mathbb{R}^d$ 的闭合 $\to$ $\mathbb{R}^n$ 的不闭合
\end{itemize}

\textbf{残差加法的功能必然性：}残差连接 $x^{(l+1)} = x^{(l)} + F(x^{(l)})$ 是一个离散Euler积分器。为什么不用乘法？若 $x^{(l+1)} = x^{(l)} \odot F(x^{(l)})$（乘法残差）：$|F| < 1$ 时轨线指数衰减至0（深度坍缩），$|F| > 1$ 时轨线指数发散（表示爆炸）。加法残差通过 Jacobian $= I + \partial F/\partial x$ 使特征值维持在 $\approx 1$ 附近——\textbf{加法是唯一能长期维持复杂轨线的累积方式}。这一事实将在第\ref{sec:edge_chaos}节（边缘混沌）中获得完整的动力学解释。

\textbf{对"弱闭合"的重新定位：}传统上残差加法与态射乘法的"不匹配"被描述为框架的最薄弱环节。通过本节的分析，这一不匹配获得了精确的结构刻画：它是 $\mathbb{R}^n \otimes \mathbb{R}^d$ 分解中不对称闭合的直接表现，且此不对称性不是需要修复的缺陷——它是Transformer能够产生语义动力学的数学前提。如果两端都闭合（A可复合）$\to$ 语义死亡；如果两端都不闭合（W不可复合）$\to$ 无语义自足结构 $\to$ 深度坍缩为单层。只有 $\mathbb{R}^d$ 闭合而 $\mathbb{R}^n$ 不闭合——即实际架构——才能同时拥有自足的语义空间和在演化中的语义空间上动态决策的路由系统。

\textbf{两项独立研究的经验支撑：}Anthropic的Transformer Circuits框架 \cite{elhage2022transformer_circuits} 的路径展开定理——将整个模型的前向计算精确表示为所有可能路径上的矩阵乘积之和——从机械可解释性角度证明了残差流的线性性质：非相邻层之间通过 $W_I^{\text{后}} \cdot W_O^{\text{前}}$ 构成隐式"虚拟权重"，使得加性累积在数学上诱导出乘性复合效应。Deep Delta Learning \cite{zhang2025deep_delta} 从反向确认了纯加性残差的结构性局限：标准残差连接"没有直接机制来替换或删除已过时的内容"，而通过引入门控参数 $\beta$ 统一加性跳跃与乘性覆写，可获得可量化的性能增益。两项工作从"加性何以能够工作"和"加性确实有限"两个方向为不对称闭合的立场提供了独立的经验支撑。

\textbf{论证强度：}严格层（因果依赖是数学事实，残差加法的功能分析是数学推导）+ 形式论证层（不对称闭合作为功能条件的论证）。

% -------------------------------------------------------------------
\subsection{精简原则的概率结构证明}
\label{sec:precision_proof}

第\ref{sec:intro}节的P1和P3分别陈述了精确优先于简短（精）和文本最小化（简）的使用原则。本节从概率结构出发，给出这两条原则的统一证明。

\subsubsection{错误累积：逐token的概率链}

自回归生成的过程可建模为一条概率链。设提示词经分词得到token序列 $t_1, t_2, \ldots, t_n$。模型逐token计算 $P(t_{k+1} \mid t_1, \ldots, t_k)$。在任何一步 $k$，模型以概率 $\varepsilon_k$ 产生一个偏离预期语义轨迹的预测——$\varepsilon_k$ 不为零，因为条件分布总有尾部，且训练收敛的覆盖范围有限（见第\ref{sec:truncated_projection}节）。经过 $n$ 步后，整个序列的语义仍保持在预期轨迹内的概率上界为：

\begin{equation}
P(\text{轨迹保持}) \leq \prod_{k=1}^{n} (1-\varepsilon_k)
\end{equation}

该乘积随 $n$ 的增加而单调递减。这不是说"长提示必然导致错误"——而是说，每增加一个token，就增加了一个潜在的偏离节点。错误概率不因提示词的"质量"而消失，它来自条件概率分布本身的非零尾部。在概率结构层面，$n$ 本身就是一个风险因子：较短的提示词经过较少的条件节点，累积偏离风险较低。

\subsubsection{条件坐标的覆盖稀疏性}

每一个条件上下文 $(t_1, \ldots, t_k)$ 对应概率单纯形中的一个点——模型需要在这个点上产生正确的条件分布 $P(\cdot \mid t_1, \ldots, t_k)$。然而，参数空间的有效维度仅有数百维量级（第\ref{sec:truncated_projection}节的截断投影约束）。当 $n$ 增长时，不同条件上下文的组合数以 $|V|^n$ 的速度膨胀——绝大多数可能的条件坐标在训练中从未被充分访问。这些坐标落在收敛覆盖范围之外：模型在这些位置上的条件分布不受训练收敛的约束，输出行为不可预测。

提示词越长，经过的"收敛盲区"坐标越多。这不是"模型理解力有限"——这是条件坐标空间的组合爆炸与参数空间有限维数之间的几何不匹配。\textbf{精简不是审美偏好：简（$n$ 最小化）减少累积错误节点数和收敛盲区暴露数；精（每个保留的token承载最大的信息精度）使较短序列仍能充分约束模型的行为。二者是同一个概率结构约束的两个操作面。}

% -------------------------------------------------------------------
\subsection{语义等同概率邻域：来自同音错别字的直接证据}
\label{sec:typo_evidence}

模型在处理中文时频繁输出同音错别字——"在"写成"再"、"的"写成"得"。这些字的词典定义完全不同。如果模型通过查字典式的"词义"来选字，这类错误不应该以如此高的频率出现。

错别字之所以发生，是因为在模型的训练语料中，两个同音字的条件概率分布在统计上高度重叠：它们出现在相似的句法位置、相似的前后搭配、相似的语用场景中。模型不查字典——它在每个位置上计算 $P(\text{字} \mid \text{上文})$ 并采样。当两个字的条件分布足够接近时，采样过程无法稳定地区分它们。\textbf{错别字是模型不按"词义"操作的活证据：它按概率邻域操作。}

这一现象反过来证明了正面的命题：一个语言单元在模型内部的"语义身份"，等同于所有其他单元到它的条件概率关系的总和。两个字的条件概率分布越相似，模型就越"认为"它们可以互换——不管它们的词典定义是什么。这正是分布语义学的核心命题（Firth 1957），但在本文中不需要将其作为外部理论引入——同音错别字现象本身提供了足够直接的证据。

该命题的操作含义在第\ref{sec:intro}节P2中已有体现：词汇的效力不源于其字面语义定义，而源于该词汇在训练数据中典型出现的语用场景。

% -------------------------------------------------------------------
\subsection{收敛的操作化定义}
\label{sec:convergence_def}

\begin{definition}[训练收敛]
训练收敛指模型参数在梯度下降的迭代驱动下，以训练数据上的损失函数为引导，逐步进入稳定配置的过程。操作化判据为以下三项的同时满足：(1)损失曲线停止显著下降，进入平台区域；(2)参数更新幅度趋近于零；(3)模型在训练分布内的输出行为转入可预测的稳定模式——对同分布样本的响应不再出现显著的分布漂移。
\end{definition}

\textbf{核心区分：}并非所有进入模型的信息都经历了上述收敛过程。模型参数——具体而言，$\mathbb{R}^d$ 端的W矩阵（第\ref{sec:semantic_space}节）——经历了梯度反馈循环的反复锻造；上下文窗口中的信息——$\mathbb{R}^n$ 端的路由实例（第\ref{sec:routing_space}节）——从未经历单步梯度反馈。本文中"受收敛约束"的表述，仅指相关信息经历了上述由损失函数引导的迭代稳定化过程，不附加全局最优性、唯一性或函子性等更强的数学承诺。\textbf{收敛二分的结构根源可追溯至第\ref{sec:residual_decomp}节的不对称闭合：参数} = $\mathbb{R}^d$ \textbf{端的W矩阵（固定，范畴闭合）；上下文} = $\mathbb{R}^n$ \textbf{端的路由实例（每次动态参数化，跨层不复合）。}

\begin{remark}[收敛与函子性的分离]
训练收敛保障的是稳定性（在训练分布内的一致可预测行为），而非函子性（$F(g \circ f) = F(g) \circ F(f)$）。交叉熵训练目标 $\mathcal{L} = -\sum \log P(\text{next\_token}|\text{context})$ 促使模型匹配训练数据中的条件分布——它不促使模型满足态射复合的结构保持条件。经验观察表明，在自然语言数据上收敛达到的稳定配置呈现出近似函子的结构保持性质；但这是经验发现，非收敛定义的逻辑后承。本文所有后续推导仅依赖"收敛约束赋予稳定性范围"这一性质，不使用"训练找到函子"作为任何推导的前提。
\end{remark}

% -------------------------------------------------------------------
\subsection{截断投影}
\label{sec:truncated_projection}

\begin{proposition}[截断投影约束]
参数空间的有效维度 $d_{\text{eff}}$ 在量级上显著小于其名义维度 $d_{\text{nominal}}$。模型经训练收敛后达到的稳定配置，构成完整语义空间在低维参数流形上的截断投影。超出该投影覆盖范围的语义映射，其映射行为不受训练收敛的约束。
\end{proposition}

\textbf{经验证据：}Li等人 \cite{li2018intrinsic} 测量了神经网络目标景观的本征维度，发现其约为参数总数的0.1\%至1\%。Aghajanyan等人 \cite{aghajanyan2020intrinsic} 将这一分析扩展至预训练语言模型，证明仅优化约200个参数即可达到全参数微调90\%的性能，且更大规模的预训练模型具有更低的本征维度。LoRA \cite{hu2021lora} 从工程角度验证了低秩参数化的有效性。Xu \cite{xu2026parameter} 指出学习轨迹限制在2至6维的全局子空间内。Mabrok \cite{mabrok2026rank} 进一步确认了参数空间秩的有限性。

\textbf{截断投影的结构后果：}参数空间的有限维性导致两个相互关联的结构现象。其一，收敛覆盖具有不可消除的边界——始终存在超出截断投影覆盖范围的语义映射。其二，当基座模型被用于直接对话时，输入提示极易超出截断投影的覆盖范围，由此产生的输出——在训练分布中无对应收敛约束的续写模式——构成结构性而非偶然性的系统脆弱性。人类自然语言本身亦不具有内在的逻辑结构，仅构成对逻辑的模拟；模型在语言层面的逻辑映射未经严格收敛约束，与人类语言的这一特征在结构上具有平行性。

\textbf{与第\ref{sec:residual_decomp}节分解的衔接：}截断投影描述的是 $\mathbb{R}^d$ 端语义空间的容量限制——W矩阵的低秩流形约束了可被可靠编码的语义态射范围。第\ref{sec:routing_space}节描述的是 $\mathbb{R}^n$ 端路由空间的不可累积性——即使语义空间无限维，路由本身的逐层重参数化也构成独立的可靠性限制。两者正交且互补：语义空间有限（维度约束）+ 路由空间不复合（结构约束）$\to$ 任何单次前向传播的信息处理能力受双重限制。

% -------------------------------------------------------------------
\subsection{记忆的结构性二分}
\label{sec:memory_dichotomy}

\begin{proposition}[记忆的结构性二分]
LLM Agent的记忆由两种在数学结构上根本不同的存储形式构成：
\begin{enumerate}[leftmargin=*,itemsep=0.3em]
    \item \textbf{长期记忆（参数化存储）：}参数空间中的一个点——其信息组织方式为空间结构。由 $\mathbb{R}^d$ 端的W矩阵承载。经历了训练收敛的约束。训练完成后静态冻结，不再随使用过程而演化。信息访问模式为分布式并行检索。
    \item \textbf{短期记忆（上下文存储）：}序列中的线性展开——其信息组织方式为顺序结构。由 $\mathbb{R}^n$ 端路由在语义空间 $\mathbb{R}^d$ 中产生的轨线实现。从未经历梯度反馈的锻造。具有顺序性，不具有时间性。信息访问模式为序列性的注意力加权检索。
\end{enumerate}
\end{proposition}

两种存储的差异不在容量的大小，而在数学结构的异质性。前者是嵌入在高维流形中的静态空间点，后者是指定长度上的线性符号序列。二者属于不同的数学范畴。\textbf{不存在一种"中期记忆"结构构成二者之间的连续过渡。}当前文献中所有被称为"中期记忆"的工程方案——包括MemGPT的层级存储管理 \cite{memgpt_2023}、Titans的神经长期记忆模块 \cite{titans_2024} 以及HAMR框架 \cite{hamr2024}——本质上是在管理短期记忆的容量分配，即在上下文窗口的内外之间执行信息的换入换出操作，而非在参数化空间结构与序列化顺序结构之间创建第三种具有独立数学性质的存储类别。

\textbf{时间性维度的缺失：}Transformer中的位置编码提供的是顺序信息（token A位于token B之前——一个离散的二元关系），而非时间性信息（token A距离当前时刻有多远——一个连续的、具有心理物理含义的量，以及该距离如何影响信息的可及性、可信度和关联强度）。在注意力机制下，所有位置的token对任意查询均平等可及——不存在基于时间距离的自发衰减、不存在基于重要性的自发巩固、不存在基于经验的自发重组。这与人类记忆系统形成结构性对比：人类记忆中信息的衰减曲线、睡眠依赖的巩固过程和基于相似性的联想重组由生物记忆系统的底层机制自动执行，无需刻意管理。

\begin{remark}[精确范围的限定]
上述"无时间性"论断针对的是跨会话的时间性动态。在单个上下文窗口内部，注意力机制确实表现出特定的顺序性效应——时间注意力头（temporal heads）对近期token的偏好、近因效应（recency bias）等 \cite{liu2023lost}。这些窗口内的顺序性效应构成上述论断的局部限定，但其有效范围限于单窗口内的顺序性偏好，不扩展至跨窗口的时间性连续记忆。
\end{remark}

% -------------------------------------------------------------------
\subsection{边缘混沌：不复合的动力学后果}
\label{sec:edge_chaos}

第\ref{sec:routing_space}节确立了注意力跨层不复合是数学事实。本节分析这个事实的动力学后果——不复合如何产生Transformer的复杂行为，以及训练良好的Transformer处于何种动力学状态。

\subsubsection{不复合驱动语义聚类}

注意力不复合意味着路由每层重新决策，不被"平均化"约束。这产生了一个正反馈环：

\begin{quote}
相似token $\to$ 相似Q/K向量 $\to$ 相似注意力模式 $\to$ 从上下文收取相似信息 \\
$\to$ 相似更新 $\to$ 保持相似 \quad （\textbf{聚类自持}——收敛面）

相异token $\to$ 不同注意力模式 $\to$ 不同更新 $\to$ 更加相异 \quad （\textbf{聚类分化}——扩张面）
\end{quote}

如果注意力可复合（Markov链），长程将收敛到平稳分布，所有token的注意力行相同，正反馈被平均化消灭。\textbf{不复合 = 聚类可以发生的条件。}

\subsubsection{混沌通路的三个条件}

不复合提供了\textbf{路径依赖}（第一步往哪走影响第二步往哪走——路由不由固定转移概率决定，而由当前语义状态决定）。但路径依赖单独不足以产生混沌——还需要两个条件，而Transformer恰好都满足：

\begin{table}[H]
\centering
\caption{混沌三条件在Transformer中的实现}
\begin{tabular}{@{}p{2.5cm}p{5cm}p{5cm}@{}}
\toprule
\textbf{条件} & \textbf{数学机制} & \textbf{Transformer中的实现} \\
\midrule
路径依赖 & 下一步依赖当前状态 & $\mathbb{R}^n$ 不米田 $\to$ $a_{ij}$ 依赖 $\mathbb{R}^d$ 当前状态 \\
非线性折叠 & 微小差异被放大为大位移 & Softmax将 $Q \cdot K$ 小差异放大为概率单纯形大位移；ReLU/GeLU阈值开关；LayerNorm约束模长到球面 \\
局部扩张 & Jacobian在某方向特征值 $> 1$ & 残差Jacobian $J = I + \partial F/\partial x$，在某些方向上特征值 $> 1$，相邻轨线指数分离 \\
\bottomrule
\end{tabular}
\end{table}

三者同时成立 $\to$ \textbf{混沌通路}：微小的初始语义差异在通路迭代中被放大，相似起点的token在短期共向（收敛面），长期分叉（扩张面）。\textbf{不米田是混沌的必要条件，不是充分条件}——如果F是线性的，路径依赖的系统只会收敛或发散，不会混沌。

\subsubsection{边缘混沌假说}

将残差流视为离散动力系统：

\begin{equation}
x^{(l+1)} = x^{(l)} + F(x^{(l)}; W)
\end{equation}
\begin{equation}
J^{(l)} = I + \frac{\partial F}{\partial x}\Big|_{x=x^{(l)}}, \quad \lambda_i = \lim_{L \to \infty} \frac{1}{L} \sum_{l=1}^L \ln \sigma_i(J^{(l)})
\end{equation}

\begin{table}[H]
\centering
\caption{Lyapunov谱与Transformer动力学状态}
\begin{tabular}{@{}p{2.5cm}p{4.5cm}p{5cm}@{}}
\toprule
\textbf{动力学状态} & \textbf{Lyapunov谱} & \textbf{对Transformer意味着什么} \\
\midrule
$\lambda_{\max} \ll 0$ & 过阻尼 & 轨线收敛到不动点 $\to$ 深度坍缩 $\to$ 所有token最终同质 \\
$\lambda_{\max} \approx 0$ & 边缘混沌 & 轨线既不收敛也不发散 $\to$ 深度有持续信息增量 $\to$ 语义结构可维持 \\
$\lambda_{\max} \gg 0$ & 混沌 & 轨线指数发散 $\to$ 表示爆炸 $\to$ 初值极端敏感 $\to$ 不可训练 \\
\bottomrule
\end{tabular}
\end{table}

\textbf{训练 = 在参数空间中寻找边缘混沌态。}训练好的Transformer必须恰好坐在 $\lambda_{\max} \approx 0$ 的临界线上。

这解释了已知但未系统解释的现象：
\begin{itemize}[leftmargin=*,itemsep=0.2em]
    \item \textbf{Over-depth退化：}层数超过Lyapunov时间后，轨线在混沌吸引子上随机游走，新层不增加信息
    \item \textbf{残差不能去掉：}去掉恒等项I $\to$ Jacobian失去对角优势 $\to$ 谱失控 $\to$ 无法维持在边缘混沌
    \item \textbf{LayerNorm的位置敏感性：}LayerNorm = 每步对状态做投影 $\to$ 约束轨线不跑出有界区域 $\to$ 混沌保持"有界非周期"
    \item \textbf{越深越难训：}长轨线对参数极度敏感 $\to$ 反向传播通过边缘混沌系统 $\to$ 梯度本身也边缘混沌
    \item \textbf{残差加法而非乘法的动力学必然性：}乘法残差无法维持 $\lambda_{\max} \approx 0$——$|F|<1$时所有轨线指数衰减至0，$|F|>1$时所有轨线指数发散。只有加法残差（$J = I + \partial F/\partial x$）允许特征值维持在 $\approx 1$ 附近，使边缘混沌可持续
\end{itemize}

\textbf{论证强度：概念构造层。}路径依赖和不复合的推导属于严格层；非线性折叠和局部扩张的识别属于形式论证层；Lyapunov谱的分析是概念推导——三条件的"同时成立"判断来自对Transformer计算图的逐项数学分析而非数值验证。整体标注为概念构造层，但给出了可检验的预测（T7--T10，第\ref{sec:predictions}节）。

% -------------------------------------------------------------------
\subsection{使用原则的统一推导}
\label{sec:principles_derivation}

第\ref{sec:intro}节陈述的十七条使用原则在本节的各项基础事实上获得了统一的证明。以下逐条给出推导要旨，每条仅标注其所依赖的基础事实。

\medskip
\textbf{第一组：输入原则（P1--P3）}

\textbf{P1（精确优先于简短）：}已证于第\ref{sec:precision_proof}节。错误累积 $\prod_{k=1}^n (1-\varepsilon_k)$ 随 $n$ 单调递减；条件坐标以 $|V|^n$ 膨胀而有效维度仅数百。故 $n$ 最小化降低累积偏离风险，每个保留的token须承载最大信息精度——精与简是同一概率约束的两个操作面。

\textbf{P2（词汇的分布场效应）：}已证于第\ref{sec:typo_evidence}节。同音错别字显示模型按 $P(\text{字} \mid \text{上文})$ 采样而非按词典定义选字。故一个词的语义身份等同于所有其他词到它的条件概率关系的总和。

\textbf{P3（目标约束而非路径约束）：}每增加 $m$ 个路径指令token，错误累积从 $\prod^n$ 放大为 $\prod^{n+m}$（事实1），且路径指令中每个词汇自带概率邻域（P2——事实1的直接推论），邻域叠加增加语义漂移风险。路径指令增加条件节点数和歧义密度；目标指令仅锚定输出端——不增加条件链上的偏离节点。

\medskip
\textbf{第二组：上下文管理原则（P4--P5）}

\textbf{P4（上下文信息的审慎选择）：}上下文信息未经历梯度筛选（事实2：收敛二分——$\mathbb{R}^d$参数闭合而$\mathbb{R}^n$路由不复合），注意力对所有窗口内token平等可及（事实1）。冗余信息与关键信息在softmax分母中无优先级区分地竞争注意力份额。故注入的每条信息须经审慎评估：选错的代价不是浪费空间，是稀释关键约束。

\textbf{P5（上下文的分级滚动管理）：}P4在动态场景中的推广。窗口容量固定且注意力不加区分$\to$在有限容量下最大化关键信息可及性的严格策略为优先级分级、高常驻低滚动。

\medskip
\textbf{第三组：收敛任务原则（P6--P10）}

\textbf{P6（能力边界的节点划分）：}$\mathcal{I}(T) > \effcap$ 时部分约束超出窗口表示范围$\to$这些维度上的输出不受约束。分解使每个 $\mathcal{I}(T_i) \leq \effcap$，各约束重新进入注意力可及范围（事实1）和收敛覆盖范围（事实3）。这是恢复收敛工作条件的操作，非性能优化。

\textbf{P7（测试收敛驱动）：}参数冻结=训练收敛通道关闭（事实2）。Agent运行时唯一残留的修正机制为生成$\to$独立验证$\to$反馈$\to$再生成循环。无验证信号时此循环退化为无导向采样。

\textbf{P8（外部约束锚定）：}分解决策本身在上下文中做出，上下文信息未经收敛（事实2）$\to$分解决策不可靠。其结构根因在 $\mathbb{R}^n$ 跨层不复合（第\ref{sec:routing_space}节）：路由决策不经收敛约束，在此基础上的分解决策同样不经收敛约束。父节点偏差 $\delta$ 沿层级乘性放大（事实3：截断投影边界外的语义映射无约束$\to$层级间信息损失每层累积）。外部锚定（人工判断或经独立测试的程序逻辑）是唯一可切断此偏差传导链的信息源。

\textbf{P9（生成与检验的分离）：}检验是上下文中的推理，未经收敛（事实2）。若生成与检验共享上下文，模型在生成中犯的错误在检验中同样的条件分布下大概率重现——因两次推理之间无梯度信息修正该条件分布。分离检验的上下文/模型实例则检验信息集独立于生成过程的条件依赖，盲区被打破。

\textbf{P10（逐环节验证）：}P8的时序操作化。每层产出后立即验证修正$\to$ $\delta_k$ 不进入 $\delta_{k+1}$ 计算。逐环节验证不降低每层 $\varepsilon$；它阻止 $\varepsilon$ 跨层乘积放大（事实3的乘性传导）。

\medskip
\textbf{第四组：发散任务原则（P11--P12）}

\textbf{P11（发散与收敛的空间隔离）：}偏好信号（相对序）与验证信号（绝对坐标）信息结构不同。共享上下文时二者在注意力计算中交叉污染（事实1：注意力不加区分地对所有token加权）。分离确保两种异质信号各自在其适配的注意力空间内运作。

\textbf{P12（偏好驱动的对话式协作）：}$\verifspace = \emptyset$ 时测试收敛失去工作条件（P7的推论）。偏好信号——操作者对候选方向的比较性判断——是此时唯一残留的导向源。操作者在线性是收敛循环闭合条件的直接推论：验证信号缺失时，人的偏好信号替代验证信号闭合迭代环。

\medskip
\textbf{第五组：结构边界原则（P13--P17）}

\textbf{P13（扁平化管理）：}子Agent的驱动提示词是上层Agent的上下文推理产物——未经收敛（事实2），不可靠性为 $\varepsilon_d$。$k$ 层嵌套后顶层驱动提示词的不可靠性累积为 $1 - \prod_{i=1}^k (1-\varepsilon_{d,i})$，且每层驱动消耗上层窗口的注意力容量。扁平化将 $k$ 约束为1。

\textbf{P14（操作者上限）：}Agent的不可自举节点——根需求锚定（P16）、关键分叉仲裁（P8）、偏好信号注入（P12）——均由操作者填补。操作者不能传递超出自身认知能力的信息。故系统上限=操作者在这些节点处的认知质量上界。

\textbf{P15（验证覆盖即能力边界）：}测试收敛有效性条件为 $\verifspace = \taskcspace$（事实3的直接推论：未被 $\verifspace$ 覆盖的约束维度无收敛约束$\to$Agent可系统性违反）。故Agent可靠性的结构上界为 $|\verifspace|/|\taskcspace|$。这是信息论恒等式。

\textbf{P16（根需求的外在性）：}验证链顶端的判断——"当前产出是否符合用户真实意图"——所依赖的语义组合（该具体用户在该具体时刻的全部语义约束）在训练数据的概率分布中不具有系统性标注。模型对该判断缺乏梯度信息（事实2的推论：未经损失函数引导锻造的语义态射不受收敛约束）。故验证链在此处不可闭合。

\textbf{P17（幻觉边界不可消除）：}softmax在任何输入上必然输出一个概率分布（事实1：它不会输出"我不知道"）。但 $d_{\text{eff}}$ 始终有限（事实3）$\to$ 收敛覆盖始终有边界$\to$对边界外的输入，模型产生的输出不受收敛约束。更好的训练可扩大覆盖面积，但不可消除边界本身。

\medskip
\textbf{总结：}十七条原则仅依赖五个基础事实。其中事实1--4为严格数学事实或形式论证，事实5（不对称闭合——第\ref{sec:causal_coupling}节）为已有精确结构刻画、功能含义已澄清的环节。所有推导遵循"基础事实$\to$原则$\to$操作含义"的单向链，无循环依赖。


% ===================================================================
\section{理论与猜想定义}
\label{sec:definitions}

\subsection{核心概念的形式定义}

在展开形式推导之前，本节对全文的核心概念给出精确的操作化定义。这些定义的目的在于将前文中的直觉用语替换为具有明确外延和操作化判据的技术概念。

\begin{table}[H]
\centering
\caption{核心概念定义表}
\begin{tabular}{@{}p{2.5cm}p{5.5cm}p{4cm}@{}}
\toprule
\textbf{概念} & \textbf{形式定义} & \textbf{来源} \\
\midrule
残差流 & Transformer层状态 $X^{(l)} \in \mathbb{R}^n \otimes \mathbb{R}^d$，携带语义（$\mathbb{R}^d$）和路由（$\mathbb{R}^n$）两种正交且因果耦合的信息 & \S\ref{sec:residual_decomp} \\
语义空间（$\mathbb{R}^d$） & 残差流的 $\mathbb{R}^d$ 因子——token含义的向量空间，W矩阵在其上构成范畴（复合闭合，米田成立） & \S\ref{sec:semantic_space} \\
路由空间（$\mathbb{R}^n$） & 残差流的 $\mathbb{R}^n$ 因子——token间注意力路由的空间，A矩阵跨层不构成范畴（复合失败，米田不适用） & \S\ref{sec:routing_space} \\
不对称闭合 & $\mathbb{R}^d$ 范畴闭合 + $\mathbb{R}^n$ 范畴不闭合 + 因果耦合（$\mathbb{R}^d$ 闭合导致 $\mathbb{R}^n$ 不闭合） & \S\ref{sec:causal_coupling} \\
训练收敛 & 参数经梯度下降在训练分布上达到损失稳定、更新趋零、行为可预测的配置 & \S\ref{sec:convergence_def} \\
测试收敛 & 上下文窗口内的"生成$\to$验证$\to$分析$\to$修正"闭环迭代，由可自动执行的验证操作驱动 & \S\ref{sec:formal} \\
可验证目标 & 可被自动执行的操作或判据判定为"已达成"的任务完成标准；为测试收敛的必要条件 & \S\ref{sec:formal}.4 \\
收敛约束 & 参数经历梯度反馈锻造后享有的结构稳定性（操作化判据：损失平台、更新趋零、行为稳定） & \S\ref{sec:convergence_def} \\
上下文窗口 & Agent在单次推理中注入前向计算的序列化信息空间，其内容不经梯度筛选 & \S\ref{sec:convergence_def} \\
有效容量 $\effcap$ & 模型在单次前向传播中实际可可靠利用的上下文长度，显著小于名义窗口上限 & \S\ref{sec:formal} \\
收敛覆盖 & 截断投影覆盖范围内的语义映射集合；该范围内的映射受训练收敛约束 & \S\ref{sec:truncated_projection} \\
分布外暴露风险 & 输入超出收敛覆盖范围时输出不受收敛约束的现象——softmax仍分配概率但无"正确/错误"的可操作判据 & \S\ref{sec:truncated_projection} \\
语义聚类 & 相似token在 $\mathbb{R}^d$ 中相互靠近、相异token相互远离的正反馈结构——由 $\mathbb{R}^n$ 不复合驱动的动力学涌现 & \S\ref{sec:edge_chaos} \\
边缘混沌 & $\lambda_{\max} \approx 0$ 的动力学状态——训练好的Transformer所处的不收敛也不发散的临界动力学区域 & \S\ref{sec:edge_chaos} \\
收敛任务 & 目标可被一组有限、可自动执行的验证标准完整刻画的Agent任务 & \S\ref{sec:formal} \\
发散任务 & 无可自动执行验证标准的Agent任务；目标在探索过程中逐渐成形，由比较性偏好驱动 & \S\ref{sec:divergent_formal} \\
根需求 & 验证链递归上溯的终点——"当前产出是否符合用户真实意图"——该语义位于训练分布之外 & \S\ref{sec:formal}.6 \\
关键分叉点 & 任务结构中不可逆的决策节点（需求消歧、接口定义、模块边界），一旦确定则锁定全部下级路径 & \S\ref{sec:formal}.3 \\
偏好信号 & 发散任务中操作者对候选方向的比较性判断（形式为"A优于B"），功能等价于采样器的方向梯度 & \S\ref{sec:divergent_formal} \\
级联传导 & 父节点语义偏差一次性向下传播至所有子节点，子节点在偏差基底上执行且缺乏独立纠正信息 & \S\ref{sec:formal}.2 \\
AGI（本文用法） & 具备全模态完整语义空间、自发求知动机、以及时间性渐进记忆结构的人类水平通用认知主体 & \S\ref{sec:agi} \\
\bottomrule
\end{tabular}
\end{table}

\subsection{框架猜想的形式陈述}

\textbf{C1（收敛二分猜想，框架性主张）：}
参数经历训练收敛而上下文未经收敛——这一二分构成对Agent行为差异的统一解释框架。结构根源追溯至第\ref{sec:residual_decomp}节的不对称闭合：$\mathbb{R}^d$端W矩阵闭合（参数），$\mathbb{R}^n$端路由不复合（上下文）。此为框架性主张而非单一可检验猜想；其综合有效性由C2至C8各项猜想的检验结果联合支撑或否证。

\textbf{C2（截断投影猜想）：}
Transformer参数空间的有效维度限制使得模型仅能学到完整语义空间在低维参数流形上的截断投影。超出投影覆盖范围的语义映射，其映射行为不受训练收敛约束。可检验推论：模型性能在训练分布的边缘区域出现系统性衰减，衰减速率与输入偏离训练分布中心的程度呈正相关。

\textbf{C3（根需求外在猜想）：}
在参数冻结、截断投影限制和上下文不经收敛三项联合前提下，全自动闭环在根需求节点处不可闭合——最高层验证标准必须由系统外部提供。否证条件：在开放域任务中发现无需人工提供根需求即可长期保持人类水平可靠性的全自动闭环系统。

\textbf{C4（发散任务偏好依赖猜想）：}
发散任务中，偏好信号的密度与方向一致性构成探索效率的充分统计量。稀疏偏好导致条件分布更新不足，探索停滞于局部区域；矛盾（不一致）偏好导致条件分布在多个不兼容梯度方向间震荡，产出发生系统性漂移；指向训练分布外区域的偏好导致模型无对应的受收敛约束的映射，输出退化。可经由受控实验检验（\S\ref{sec:empirical}.2，T2）。

\textbf{C5（自主性-可靠性互斥猜想）：}
Agent自主性程度的增加等价于将更多决策节点从外部约束转移至上下文中的测试收敛执行，等价于系统可靠性下界的系统性降低。可检验部分：自主性程度与任务完成可靠性之间存在负相关关系，且该关系不由模型规模增加所消除（\S\ref{sec:empirical}.2，T6）。"不可消除"部分为框架性主张——在"消除"的操作化判据精确建立之前，框架主张该权衡仅可放宽而不可根除。

\textbf{C6（Agent Scaling先行平台期猜想，待精确化）：}
Agent在多步任务上的性能随模型规模增长的提升曲线，将系统性地早于单步预测任务的性能曲线进入平台区域。前者的上界由验证信号对任务需求语义空间的覆盖度 $\verifspace/\taskcspace$ 决定，后者由参数空间的有效维度 $d_{\text{eff}}$ 决定。当前为方向性预测，精确化路径见\S\ref{sec:open}.6。

\textbf{C7（AGI不可达猜想）：}
在当前Transformer训练范式下——参数部署后冻结、softmax注意力作为唯一的结构化信息传递机制、非语言模态通过联合分布拟合进行形式对齐——LLM Agent不可能达到人类意义上的通用人工智能。该结论建立在五项结构性缺失的独立论证链之上：(1)语义空间不完备（截断投影），(2)模态因果结构不对齐（注意力限于语言），(3)动机不自发（根因未明，仅为观察），(4)无渐进记忆（记忆结构二分），(5)顺序性而非时间性（位置编码的结构限制）。v4新增第(6)项——\textbf{路由不复合}：$\mathbb{R}^n$ 跨层不构成范畴 $\to$ 路由决策无累积 $\to$ 无稳定行为策略的跨情境迁移（\S\ref{sec:routing_space}）。本猜想仅针对当前范式，不声称在任何可能架构下AGI均不可达。

\textbf{C8（边缘混沌猜想，v4新增）：}
训练良好的Transformer处于边缘混沌态（$\lambda_{\max} \approx 0$）。过阻尼（$\lambda_{\max} \ll 0$）导致深度坍缩和表示同质化；混沌（$\lambda_{\max} \gg 0$）导致表示爆炸和不可训练。此猜想为C1--C2的动力学基础提供统一解释，并通过第\ref{sec:routing_space}节的不复合论证与第\ref{sec:edge_chaos}节的混沌三条件分析建立与框架数学基础的连接。可检验预测见\S\ref{sec:predictions}（T7--T10）。


% ===================================================================
\section{证明与论证}
\label{sec:formal}

\subsection{形式推导的前提系统}

以下前提构成了收敛任务理论的全部公理基础。每条前提均标注了其来源与性质。

\begin{table}[H]
\centering
\caption{形式推导前提（Premises）}
\begin{tabular}{@{}c@{}p{7.5cm}@{}p{3.5cm}@{}}
\toprule
\textbf{编号} & \textbf{陈述} & \textbf{来源} \\
\midrule
P0 & 模型处理任何任务的"统筹"与"连贯"——跨信息片段的关联与整合——均以有限的文字输入为基础 & 经验观察 \\
P1 & 文字输入的有效边界划定了模型对单一问题的能力边界 & \S\ref{sec:truncated_projection}的推论 \\
P2 & 超出有效边界的信息不受训练收敛约束的覆盖 & \S\ref{sec:truncated_projection} \\
P3 & Agent在运行中通过测试收敛逼近目标——可验证目标驱动反复迭代 & \S\ref{sec:definitions} \\
P4 & 测试收敛的有效运作以可验证目标的存在为前提 & \S\ref{sec:definitions} \\
P5 & 上下文窗口中的映射未经训练收敛约束，不具备结构稳定性。\textbf{结构根因：}$\mathbb{R}^n$ 跨层不复合（\S\ref{sec:routing_space}）——路由决策每层在当前语义状态下重新参数化，不享有跨层的结构稳定性 & \S\ref{sec:convergence_def} + \S\ref{sec:residual_decomp} \\
P6 & Agent的功能架构可分解为模型与Harness的组合 & 工程定义 \\
\bottomrule
\end{tabular}
\end{table}

\textbf{从P0至P1的推导：}模型赖以执行任何认知操作的唯一信息载体为文字输入序列（P0）。输入文本经由注意力映射进入参数空间；参数空间的有效维度 $d_{\text{eff}}$ 在量级上受到截断投影限制（\S\ref{sec:truncated_projection}），由此模型仅能可靠地处理信息量不超出截断投影覆盖范围的输入。文字输入的有效边界在此处被划定（P1）。P0作为P1的元前提，声明的不仅是单一问题的边界性，而是模型所依赖的媒介本身——文字序列——的结构性边界。

\subsection{桥接前提}

\begin{table}[H]
\centering
\caption{桥接前提（Bridge Premises）}
\begin{tabular}{@{}c@{}p{7cm}@{}p{3cm}@{}}
\toprule
\textbf{编号} & \textbf{陈述} & \textbf{性质} \\
\midrule
B1 & 大多数具有工程意义的任务，其完整约束描述所需的信息量超出单次文字输入的有效边界；分解因而是操作上的必然 & 经验观察 \\
B2 & 层级任务结构中，父节点的输出（分解方案、接口定义、需求消歧决定）向下一次性传导至所有子节点；父节点偏差使全部子节点在偏差基底上执行，修正成本随层级深度指数上升 & 结构推理 \\
B3 & 验证标准的递归上溯终止于一个不可由系统自身生成的权威——验证链的顶层判断需要全局语义信息，该信息位于训练分布之外 & 推导自P1+P2+P4+\S\ref{sec:truncated_projection} \\
\bottomrule
\end{tabular}
\end{table}

\textbf{B3的推导：}验证链中每一层级的验证操作在其各自的上下文窗口内执行（由P1）且仅持有该层级的局部信息（由P2）。位于验证链顶端的最终判断——"当前产出是否满足用户的真实意图"——需要综合用户在具体时刻、具体场景下的全部语义约束，而该语义组合位于模型训练分布之外（由\S\ref{sec:truncated_projection}截断投影边界）。验证链因此终止于一个不可由系统自身生成的权威源。

\subsection{推导一：任务必须分解}

\textbf{前提：}B1，P1，P2。

\textbf{推导：}具有工程意义的大多数任务，其完整约束描述的信息量超出单次文字输入的有效边界 $\effcap$（B1）。超出 $\effcap$ 的信息不受训练收敛约束的覆盖（P1 + P2）。因此，此类任务必须在 $\effcap$ 边界处分解为信息量落入边界之内的子任务。不执行分解意味着任务的某些约束维度必然处于收敛覆盖范围之外——此时在这些维度上Agent的行为不受收敛约束，输出可靠性无法得到结构性保证。

\subsection{推导二：子任务须可验证，分解自然形成层级}

\textbf{前提：}P1，P3，P4。

\textbf{推导：}子任务的目标必须在 $\effcap$ 之内设定（P1的逐个子任务实例化）。子任务的执行需要可验证目标以触发测试收敛机制（P3 + P4）。分解行为本身构成一个任务实例——它需要确定拆分方案、接口定义和验证标准——因而受同一有效边界约束，同样需要测试收敛。由此，\textbf{层级结构自然形成：上层收敛过程产出的分解方案和验证标准，构成下层各子任务的收敛目标}。

\subsection{推导三：层级节点须由外部硬约束锚定}

\textbf{大前提的独立论证：}经收敛约束的信息在任务相关分布上具有一致且可预测的输入-输出行为（\S\ref{sec:convergence_def}）。工程语境中的"可靠决策"可操作化定义为：在给定的输入变差范围内，决策输出的偏差幅度具有确定的、可预测的上界。训练收敛是使得映射行为获得此类输入-输出稳定性的唯一已知系统性机制——不经收敛的信息（上下文中的推理）在输入发生微小扰动时可能产生不可预测的输出翻转，这是截断投影（\S\ref{sec:truncated_projection}）的直接后承：位于投影边缘或之外的语义映射不具有受约束的映射行为。\textbf{其更深层的结构根因在第\ref{sec:routing_space}节：上下文中的分解决策依赖} $\mathbb{R}^n$ \textbf{端的路由选择，而路由每层重新参数化、跨层不复合——路由本身不受收敛保证，在路由基础上的分解决策亦然。}因此，一个需要稳定结构保持映射的决策，其所依据的信息基础必须经受过收敛约束的锻造。

\textbf{大前提：}需要可靠结构保持映射的决策，其信息基础必须为经收敛约束的信息。
\textbf{小前提：}任务分解的元决策——节点以何种粒度切分、子任务接口以何种契约定义、依赖关系以何种顺序编排——需要可靠的结构保持映射来保证分解方案在被分解对象的语义结构上是保真的；然而，在上下文窗口中可用于该元决策的信息未经收敛约束（P5）。
\textbf{结论：}由上下文驱动的任务分解元决策在结构上不可靠。

\textbf{级联代价的叠加（B2）：}上下文驱动的分解决策不仅不可靠，且其错误代价沿层级乘性放大——父节点偏差一次性传导至全部子节点，各子节点在偏差基底上执行且其局部上下文缺乏独立识别和纠正父节点偏差所需的信息。两项因素的联合——不可靠性（来自未经收敛的信息基础）与高代价（来自级联传导）——推出\textbf{关键分叉点必须由独立于模型上下文的外部硬约束锚定}：或为人工定义的确定性规则，或为经独立验证的程序逻辑。该结论不是从特定工程案例中归纳的经验教训，而是从"上下文不经收敛"与"级联代价"两个前提中演绎出的结构必然。

\textbf{工程实例：}在本文作者构建的长篇网文多Agent创作框架中，角色状态、故事时间线与伏笔库由独立的程序逻辑进行确定性追踪，而非依赖模型在上下文中的自行记忆与判断。关键分叉点——人物出场时机、时间跳跃位置、伏笔回收节点——全部由程序标记和校验，模型仅在其约束空间内执行生成。该实例在第\ref{sec:case}节中详述。

\subsection{推导四：根需求的人工锚定}

\textbf{前提：}P4，B3。

\textbf{推导：}测试收敛机制的有效运作以可验证目标的存在为前提（P4）。验证标准沿任务层级递归上溯——每一层级的验证标准由上一层级的收敛产出来定义（P4在层级上的自反应用）。该递归链的顶层——验证"最高层目标是否被满足"——依赖于一个不可由系统自身生成的权威源（B3）。因此，\textbf{递归终止于由操作者（系统外部主体）锚定的最高层需求。全自动闭环在根需求节点处原则上不可闭合}。

\subsection{综合：多级需求收敛结构}

综合上述四项推导，收敛任务的行为结构可总结如下：模型的单任务能力边界由文字输入的有效边界划定。现实任务通常超出该边界，因此必须在边界处执行分解。分解过程本身受同一边界约束，形成层级收敛结构：上层收敛产出下层目标，硬节点由外部约束锚定，根需求由操作者提供。层级每加深一层，收敛成本——维护层级间语义保真度所需的测试收敛循环总量——呈指数增长，而高层的综合与验证操作仍然受上下文窗口容量限制。

\begin{table}[H]
\centering
\caption{"全自动闭环不可闭合"论断的三层论证结构}
\begin{tabular}{@{}p{2.2cm}p{3.5cm}p{4.5cm}p{2cm}@{}}
\toprule
\textbf{层面} & \textbf{定义} & \textbf{不可闭合的性质} & \textbf{论证强度} \\
\midrule
目标设定层 & 系统能否自主生成最高任务目标 & 逻辑必然——根需求语义位于训练分布之外，无从收敛 & 强 \\
任务分解层 & 系统能否在给定目标下自主规划子任务结构 & 结构不可靠——上下文驱动的分解决策不受收敛约束（结构根因：$\mathbb{R}^n$ 跨层不复合） & 中 \\
执行闭环层 & 系统能否在给定验证标准下自主完成子任务 & 条件可行——若验证信号充分覆盖任务需求的语义空间，则无需人工介入 & 弱 \\
\bottomrule
\end{tabular}
\end{table}

本文的核心论证支持\textbf{目标设定层}的不可闭合性。"全自动闭环不可闭合"这一论断的精确范围限定于此层——本文不声称执行闭环层在任何条件下均不可自动化。在验证信号充分覆盖任务需求语义空间的条件下（如通过完整测试套件的代码生成），执行闭环层可在无人工介入的情况下闭合。

\subsection{从Agent的结构性缺失推导人的角色}

综合前述分析，Agent相对于人类独立任务执行者存在三项结构性缺失：

\begin{enumerate}[leftmargin=*,itemsep=0.3em]
    \item \textbf{收敛映射维度的有限性}（\S\ref{sec:truncated_projection}）：超出截断投影的语义映射不受收敛约束，模型在这些区域中的输出不具有结构保证。
    \item \textbf{记忆的时间性缺失}（\S\ref{sec:memory_dichotomy}）：记忆缺乏跨会话的自发衰减、巩固与联想重组机制，不具备时间深度。
    \item \textbf{非语言模态的因果结构缺失}（\S\ref{sec:divergent_formal}）：非语言模态信号在模型中是经由联合分布拟合进行形式对齐的结果，缺乏与物理世界的因果闭环验证。
\end{enumerate}

上述缺失直接界定了操作者在人机协作中必须承担的结构性功能：

\begin{table}[H]
\centering
\caption{从结构性缺失到操作者补偿角色的映射}
\begin{tabular}{@{}p{2.5cm}p{5.5cm}p{5cm}@{}}
\toprule
\textbf{缺失} & \textbf{操作者须补偿的功能} & \textbf{论证来源} \\
\midrule
收敛维度有限 & 约束模型在分布外区域的偶发不可靠输出；将任务空间划分至上下文窗口的有效覆盖范围内 & \S\ref{sec:formal}.3--6 \\
记忆无时间性 & 在Harness层面引入外部持久化存储与跨会话信息重新注入机制，替代系统所缺乏的时间性动态 & \S\ref{sec:memory_dichotomy} \\
模态因果缺失 & 为系统注入源自人类具身经验与生活实践的动机、偏好与审美判断——"生活体验式"信息 & \S\ref{sec:divergent_formal} \\
三项共同 & 锚定验证链的根需求——该语义内容位于模型的收敛空间之外 & \S\ref{sec:formal}.6 \\
\bottomrule
\end{tabular}
\end{table}

\textbf{操作者的角色分化：}在收敛任务中，操作者承担设定根需求、仲裁关键分叉和划分任务空间的功能——操作者的认知边界直接构成系统在该任务上的能力上限（P14）。在发散任务中，操作者承担持续注入偏好信号、在模型生成与自身判断之间形成闭环迭代的功能——操作者的偏好是模型采样过程的唯一方向梯度源（P12）。两种角色的共同基础是：它们各自填补了Agent架构中的一个结构性缺口，这些缺口不由模型能力的量变（规模增长、数据扩充）而消除。

\textbf{人机协作的非过渡性：}上述三项结构性缺失的来源是当前Transformer架构的底层数学结构——参数的部署后冻结、softmax注意力作为唯一的结构化信息传递机制、模态信号的形式对齐而非因果闭环、以及残差流 $\mathbb{R}^n \otimes \mathbb{R}^d$ 中路由空间的不复合。模型能力的量变（更大规模、更长窗口、更丰富数据）可以移动截断投影的边界位置，可以扩展收敛覆盖的面积，但不改变"边界存在"这一性质本身——正如它不改变 $\mathbb{R}^n$ 端跨层不复合的结构事实。操作者在人机协作中的构成性角色来源于这些结构性缺口的存在，而非来源于当前模型能力的"不足"——因此该角色不由模型能力的后续提升而趋于消失。

\subsection{发散任务：偏好驱动迭代}
\label{sec:divergent_formal}

\textbf{论证强度：概念构造层。}核心机制已达概念精度，整体形式化密度低于收敛框架部分。

发散任务与收敛任务的根本差异在于评价模式的缺失：任务开始时 $\verifspace = \emptyset$——不存在可自动执行的验证标准。在此条件下，测试收敛完全失去运作前提。模型的真实行为机制可以还原为纯粹的\textbf{条件概率自回归采样}：$P(\text{next\_token} \mid \text{context})$。"探索""创意""主动尝试"等拟人化描述属于描述层面的便利——它们所描述的是人在与模型的交互中所观察到的涌现效果，而非模型内部发生的计算过程。将描述语言剥离之后，模型在发散任务中执行的始终是同一基本操作：在当前上下文条件下计算下一个token的条件分布并采样。

\textbf{方向性的外部来源：}当模型处于发散模式时，其采样的方向性完全来自外部偏好信号对后续采样条件分布的连续重塑。每一次操作者的比较性判断——"方向A优于方向B"——等价于向条件分布施加了一次有向更新。这一机制与收敛任务中损失函数梯度提供方向信号在功能上平行，但在信息论层面上弱于后者：验证信号提供绝对坐标（"通过/不通过"），偏好信号仅提供相对序信息（"更好/更差"），不携带关于目标绝对位置的任何信息。

\textbf{偏好信号的信息论约束：}偏好信号的有效性受以下三项信息论约束的联合限制：
\begin{enumerate}[leftmargin=*,itemsep=0.3em]
    \item \textbf{稀疏性约束：}偏好信号的注入频率低于条件分布的自然漂移速率时，采样方向无法有效累积，探索停滞于语义空间中训练分布高密度区域附近的局部邻域。
    \item \textbf{一致性约束：}偏好信号的方向在多次判断中不一致时——前一判断指向方向A，后一判断指向与A不兼容的方向B——条件分布在矛盾梯度间震荡，产出风格和方向发生非平稳漂移。
    \item \textbf{分布覆盖约束：}偏好信号指向的方向超出模型截断投影的覆盖范围时，模型在该方向上缺乏受收敛约束的语义映射，输出进入分布外区域——即不受收敛约束的区域——质量不可预测且通常退化。
\end{enumerate}

\textbf{[留空——待补充实证：发散任务中偏好信号的密度与一致性对产出质量影响的系统性工程观察]}

\textbf{v4补注——发散任务与混沌通路的深层关联：}发散任务中偏好信号的作用与第\ref{sec:edge_chaos}节的混沌通路分析存在结构上的平行性。在收敛任务中，验证信号闭合了收敛循环——它为 $\mathbb{R}^d$ 语义空间中的轨线提供了确定性的方向校正。在发散任务中，验证信号缺失——偏好信号替代了验证信号的位置，但其信息论强度（仅提供相对序）弱于验证信号（提供绝对坐标）。此外，$\mathbb{R}^n$ 端的路由在两种任务中同样不复合——发散任务中缺乏验证信号的校正，路由漂移无法被系统性地纠正，这使得发散过程的轨线在 $\mathbb{R}^d$ 中的长期行为更接近无约束混沌扩散，而非受控的边缘混沌。发散过程中观察到的"新语义涌现"——若确实发生——可能对应于混沌轨线进入了 $\mathbb{R}^d$ 中训练分布未曾覆盖的新吸引域。但当前这一对应仅是方向性提示，尚未展开为系统论证。

\textbf{人机直觉的认知平行性：}以下观察为上述机制提供认知层面的补充说明。人类在直觉模式下的认知操作与Agent在默认采样模式下的token预测在功能层面具有可类比性：二者均依赖已建立的短路径直接映射，均为认知经济原理下的默认操作模式，均具有低认知负荷特征。但该类比严格限定于\textbf{功能层面}——人类直觉运作于具身性的因果反馈循环之中（生理唤醒、情绪标记、运动预期），而token采样缺乏对应的具身结构，故不声称二者在结构上同构。逻辑推理——无论是人类的显式推理还是模型的思维链——在两种系统中均具有高昂的认知成本（更多的计算步骤、更大的上下文消耗、更高的逆缩放风险），且均依赖外部锚定物（实验验证、编译器报错、形式符号推导、可重复观测）来获得超越系统内部自持的稳定性 \cite{kahneman2011}。

\textbf{训练-推理边界的开闭状态}是区分人类认知与Agent认知的一个关键结构性差异：人类的推理与学习过程相互交叠——每一次推理都在微调突触权重，每一次学习都发生在推理的上下文之中；Agent的训练与推理被严格分离——离线执行批量梯度下降，部署后参数静态冻结。这一差异解释了数项衍生现象：人类直觉随经验的持续累积而逐渐更新，Agent直觉冻结于训练结束时刻的权重快照；人类可将长期使用的认知工具和外部锚定物逐步内化为直觉（从"需要显式验证"过渡至"自动可靠"），Agent无法执行此类内化。操作结论：\textbf{收敛任务在根需求锚定后可在验证信号充分覆盖的条件下由Agent自主执行——验证信号闭合了收敛循环；发散任务在全程均需操作者的偏好信号在线引导——偏好信号是唯一的采样方向梯度。}

\textbf{框架限制：}发散过程中观察到的"全新语义或关联的生成"——在本文框架的早期表述中被描述为发散过程的一个特征——在当前框架中被刻画为已有语义空间的偏好引导重组，而非严格意义上的生成性。新语义如何在纯粹的条件概率采样中涌现——而非仅作为已有分布结构的重新组合被感知——构成本框架的一个开放问题（\S\ref{sec:open}.3）。

\subsection{发散向收敛过渡的结构不可靠性}

发散阶段与收敛阶段之间的过渡判断——"当前探索已足够充分，可验证目标已充分结晶，可以切换至收敛执行模式"——本身发生在上下文窗口之内，因而同样不受训练收敛约束。此判断面临两种结构上不可消除的风险：(1)过早过渡：可验证目标在偏好尚未充分结晶时被过早固定，收敛阶段执行的是不成熟的目标——收敛执行的高效率在此处转化为高代价；(2)过晚过渡：可验证目标已经充分结晶后继续发散探索，资源持续消耗而边际信息增量递减。两种风险的不可消除性与收敛任务中上下文驱动分解决策的不可靠性共享同一结构根源——做出决策所需的信息位于上下文窗口之中，而上下文信息未经收敛约束（其结构根因在第\ref{sec:routing_space}节：$\mathbb{R}^n$ 跨层不复合 $\to$ 路由决策本身不稳定 $\to$ 过渡判断同样不受收敛保证）。

过渡须经由显式的外部确认机制——操作者基于对当前产出与自身内部偏好之间距离的元认知判断，决定是否触发过渡。该判断不由模型在上下文中自主执行。

\subsection{AGI不可达的汇聚论证}
\label{sec:agi}

本节将前述各章中的分散论证链汇聚为一个针对AGI可达性的综合判断。该判断建立在六项结构性缺失的独立论证之上：

\begin{table}[H]
\centering
\caption{六项结构性缺失及其论证链（v4新增第6项）}
\begin{tabular}{@{}p{3cm}p{5cm}p{4cm}@{}}
\toprule
\textbf{缺失维度} & \textbf{论证要旨} & \textbf{论证来源} \\
\midrule
语义空间不完备 & 参数空间有限维截断导致收敛覆盖具有不可消除的边界；边界之外无结构保持映射 & \S\ref{sec:truncated_projection} \\
模态因果不对齐 & 注意力机制仅对语言模态具有天然的结构适配性；非语言模态为形式对齐结果，缺乏因果闭环验证 & \S\ref{sec:semantic_space}, \S\ref{sec:divergent_formal} \\
动机不自发 & 训练目标为条件概率匹配，参数部署后冻结；模型未表现出自发的、由内部状态驱动的求知行为——但此缺失的根因尚未被框架从数学结构中推导出来 & \S\ref{sec:open}.2 \\
无渐进记忆 & 参数化记忆与上下文记忆的结构二分排除了时间性连续过渡；不存在渐进的知识累积与经验重组 & \S\ref{sec:memory_dichotomy} \\
顺序而非时间 & 位置编码提供的是离散的顺序关系而非连续的时间性维度；注意力的平等可及性抹除了时间距离的信息论效应 & \S\ref{sec:memory_dichotomy} \\
\textbf{路由不复合（v4新增）} & $\mathbb{R}^n$ 跨层不构成范畴 $\to$ 路由决策无跨层累积 $\to$ 无稳定行为策略的跨情境迁移 & \S\ref{sec:routing_space} \\
\bottomrule
\end{tabular}
\end{table}

\textbf{[留空——待补充实证：模型在开放式、无指令引导环境中缺乏自发求知行为的系统性观察记录]}

综合上述六项论证链，在当前Transformer训练范式下——参数部署后冻结、softmax注意力作为唯一的结构化信息传递机制、非语言模态通过联合分布拟合进行形式对齐——LLM Agent不可能达到人类意义上的通用人工智能。此判断不声称在任何可能的计算架构下AGI均不可达。若未来架构在以下四个维度上做出实质性改变——(1)记忆系统获得时间性连续更新机制，(2)非语言模态信号获得因果闭环验证，(3)系统出现不依赖外部损失函数的目标自创生机制，(4)$\mathbb{R}^n$ 端路由获得跨层复合或累积能力——则上述限制需要逐项重新评估。


% ===================================================================
\section{验证与实证}
\label{sec:empirical}

\textbf{本节论证强度：经验推论层。}

\subsection{已有文献提供的独立经验证据}

以下来自独立研究的经验发现为本文的核心理论主张提供了无需依赖本文框架即可成立的独立验证：

\begin{table}[H]
\centering
\caption{独立经验证据与框架主张的对应关系}
\begin{tabular}{@{}p{4cm}p{5cm}p{4cm}@{}}
\toprule
\textbf{框架主张} & \textbf{独立经验证据} & \textbf{文献来源} \\
\midrule
参数空间有效维度远小于名义维度 & 本征维度约为参数总数的0.1\%--1\%；全局子空间限制于2--6维；LoRA以极低秩分解实现高效微调 & \cite{li2018intrinsic}; \cite{aghajanyan2020intrinsic}; \cite{hu2021lora}; \cite{xu2026parameter} \\
上下文有效利用远小于名义窗口 & "Lost in the Middle"效应——中段信息访问率系统性衰减 & \cite{liu2023lost} \\
验证信号覆盖度决定Agent可靠性上限 & 代码Agent在SWE-bench等结构化基准上系统性超越开放域Agent；数学证明$>$翻译$>$开放写作的可验证性递减 & \cite{claude_code_2025}及编码基准的广泛报告 \\
多Agent辩论提升准确性的机制 & 多Agent辩论框架独立展示了辩论对事实准确性和推理质量的提升 & 多Agent辩论研究文献 \\
测试时计算的逆缩放效应 & 更长的推理链在某些任务中反而降低性能——逆缩放是验证信号不充分时的结构预期 & \cite{gema2025inverse} \\
早期全自动Agent系统的广泛失败 & AutoGPT等完全自主模式在开放域中几乎全部未能可靠运行 & \cite{autogpt_failure_2024} \\
成熟脚手架在关键节点强制人工确认 & Claude Code在文件修改前设有人工确认门控；Cursor在关键重构处提示用户介入 & \cite{claude_code_2025} \\
\bottomrule
\end{tabular}
\end{table}

\subsection{框架的差异化可检验预测}
\label{sec:predictions}

以下预测均基于本文理论框架的逻辑演绎。若框架在某一预测上被系统性的相反证据所否证，框架的相应部分需要修正或放弃。

\subsubsection{Agent层预测（T1--T6）}

\textbf{T1（Agent任务Scaling的平台期先行）：}在控制单步预测能力增长的前提下，Agent在多步任务完成率上的Scaling曲线将系统性地早于单步预测的Scaling曲线进入增长平台期。

\textbf{T2（发散任务中偏好信号的效率决定）：}在受控实验中固定候选生成器的条件下，操控偏好反馈的频率（密度）和方向一致性，可观察到探索效率（从初始状态到满足偏好目标的平均迭代步数）随偏好密度提高而单调改善、随偏好不一致性提高而单调恶化。

\textbf{T3（发散-收敛过渡判断的可靠性下界）：}让模型在发散任务中自主判断过渡时机，随任务复杂度的增加，模型判断与人工最优判断之间的偏差将收敛至一个显著大于零的系统性下界。

\textbf{T4（发散收敛空间隔离的效应）：}采用独立上下文空间与显式过渡协议的Agent系统，在发散-收敛混合任务的完成率和产出质量上系统性优于将发散与收敛共享同一上下文空间的Agent系统。可经由A/B实验设计检验。

\textbf{T5（逆缩放与验证稀疏性的相关）：}在可按验证信号密度量化分层的任务集合中，逆缩放现象的发生率与严重程度与任务所在层级的验证信号密度呈负相关。

\textbf{T6（自主性-可靠性的结构负相关）：}Agent系统的自主性程度（操作化定义为上下文驱动决策所占比例）与任务完成可靠性之间存在系统性负相关关系，且该关系的强度不由模型规模（参数量）的增加而系统性削弱。

\subsubsection{架构底层预测（T7--T10，v4新增）}

以下四条预测测试的是框架的数学基础——不对称闭合和边缘混沌假说——而非Agent层推论。它们更底层，也更危险（更容易被否证）：

\textbf{T7（语义聚类随层增强）：}已训练Transformer各层token表示的轮廓系数（silhouette score）应随层数保持高位或递增。预测机制：$\mathbb{R}^n$ 不复合 $\to$ 正反馈（相似趋同/相异趋分）$\to$ 聚类随深度增强。若注意力跨层复合成立，轮廓系数应随层数衰减。

\textbf{T8（最大Lyapunov指数 $\approx 0$）：}训练良好的Transformer，逐层Jacobian的最大Lyapunov指数应在0附近——既不收敛（$\lambda_{\max} \ll 0$，坍缩）也不发散（$\lambda_{\max} \gg 0$，混沌）。检验方法：在标准benchmark模型上计算逐层Jacobian谱分解并估计Lyapunov指数。

\textbf{T9（残差强度调节分叉）：}将残差连接改为 $x + \varepsilon \cdot F(x)$，当 $\varepsilon \to 0$ 时恢复深度坍缩（Jacobian $\to I$，收敛到不动点，聚类消失）；当 $\varepsilon \to \infty$ 时进入混沌（Lyapunov指数 $> 0$，表示破缺）。可通过在已训练模型上插值 $\varepsilon$ 并测量表示多样性来检验。

\textbf{T10（强制跨层复合 $\to$ 同质化）：}若通过架构修改强制注意力跨层复合（如注意力参数共享 + Chapman-Kolmogorov损失），token表示的聚类结构应随深度消失，模型表现应退化。这是最直接的否证实验——如果强制复合后模型表现不退化，则"不复合 = 功能性"的核心命题（第\ref{sec:routing_space}节）被否证。

\subsection{工程案例：长篇网文多Agent创作框架}
\label{sec:case}

以下案例来自本文作者的工程实践。其性质为框架原则指导下的系统设计与演示，而非严格控制变量的因果验证实验。案例按"问题设定$\to$框架预测$\to$系统实现$\to$结果与限制"的结构展开。

\textbf{问题设定：}长篇网络小说（篇幅超过十万字量级）的AI辅助创作。该类任务的核心挑战在于：多人物线索的并行推进、跨章节伏笔的埋设与回收、以及主叙事因果链在超长文本跨度上的逻辑连贯性。在通用Agent框架（以通用对话模式或单Agent直接提示方式创作长篇）的测试中，系统性地观察到主干叙事逻辑断点——模型在累积上下文超过一定长度后丢失主线因果链的连贯性，出现前后矛盾且无法自动检测和修复。短篇与场景级写作上已有可用的Agent方案，但长篇叙事连贯性始终构成瓶颈。

\textbf{框架预测：}根据本文建立的原则——能力边界节点划分（P6）、上下文审慎注入与分级滚动（P4--P5）、外部程序硬约束（P8）、生成与检验分离（P9）、逐环节验证（P10）——一个严格遵循上述原则设计的Agent系统，应在长篇叙事的逻辑连贯性上显著优于未遵循上述原则的通用方案。

\textbf{系统实现：}
\begin{enumerate}[leftmargin=*,itemsep=0.3em]
    \item \textbf{层级节点划分（P6）：}将长篇创作任务分解为五级粒度层级——结构大纲、卷、章、节、场景。每一层级的产出须经操作者确认后方可展开至下一层级。单次生成的信息量始终被约束在模型的有效上下文容量之内。
    \item \textbf{叙事描写分通道管理（P4--P5）：}叙事主线信息（人物弧光发展轨迹、因果链依赖图谱、伏笔埋设与回收计划）在上下文中常驻于高优先级区域并随任务推进滚动更新。场景描写与对话细节按当前生成位置按需注入，使用完毕后滚动淘汰。叙事通道与描写通道的结构化隔离避免了两种异质信息在注意力计算中的相互稀释。
    \item \textbf{程序硬约束追踪（P8）：}人物状态数据库、故事时间线与伏笔注册表由独立的确定性程序逻辑进行状态管理，不由模型的上下文记忆承载。每次生成前，程序层注入当前状态的格式化快照；每次生成后，程序层对一致性约束执行自动校验。关键分叉决策——人物出场时机、时间跳跃位置、伏笔回收触发条件——由程序标记和门控，模型在约束空间内执行内容生成。
    \item \textbf{多代理生成与检验分离（P9，推论5）：}内容生成（写作）、逻辑一致性审查（检验）与文学校对（审查）分别由三个独立的Agent实例在各自独立的上下文窗口中执行。写作Agent与检验Agent不共享上下文空间——避免自写自验的循环盲区。
    \item \textbf{场景级逐环节验证（P10）：}每个场景级别的内容单元在生成后立即进入一致性校验流程，校验未通过的内容不允许流入下一场景的上下文窗口。错误在场景粒度被截断，而非累积至章节或卷的粒度后才被发现。
\end{enumerate}

\textbf{结果与限制：}在本文作者此前的工程经验基线（以通用对话模式或单Agent直接提示方式创作长篇）上，本系统运行中未再出现此前在基线条件下频繁发生的主干叙事逻辑断点。以下限制需诚实标注：(1)未执行系统化的对照消融实验——五项原则各自对总体效果的独立贡献量未被分离估计；(2)"逻辑断点"的判定基于作者对自身项目的主观评估，未引入第三方盲审或标准化的叙事连贯性量化指标；(3)该案例的性质为框架原则指导下的设计演示，而非对五项原则各自因果效应的独立验证。五项原则在该案例中以联合方式被演示——各原则之间的交互效应与每条原则的独立因果贡献需进一步的受控实验来分离和确认。

\subsection{框架的Popperian边界}

科学理论的必要条件是可被经验证据所否证。以下列出可操作的否证条件——若这些观察被可靠地确证，本文框架的相应部分将被否证或需要根本性修正：

\begin{itemize}[leftmargin=*,itemsep=0.3em]
    \item \textbf{对C3的否证：}发现一个全自动闭环系统，在开放域任务中、在无需操作者提供根需求锚定的条件下，长期且跨任务地保持人类水平的可靠性。其中"开放域"指任务空间不预先受限于一组封闭的可自动验证标准，"长期"指跨越多轮独立任务会话而性能无明显退化。
    \item \textbf{对记忆二分的否证：}一个被称为"中期记忆"的工程方案被证明使模型表现出真实的跨会话时间性动态——包括但不限于自发衰减（与时间距离单调相关的可及性降低）、自发巩固（高重要性信息随时间的访问率提升）和自发重组（基于语义相似性的跨时间信息关联）。
    \item \textbf{对C6的否证：}Agent在多步任务上的性能Scaling曲线在统计上无法与单步预测任务的Scaling曲线区分——两条曲线在全部观测的模型规模区间内保持系统性的平行，且无平台期分化的趋势。
    \item \textbf{对C8和不对称闭合核心命题的否证（v4新增）：}发现强制跨层复合的Transformer（T10）在语义聚类和下游表现上不退化——将直接否证"不复合 = 功能性"的核心命题（第\ref{sec:routing_space}节）和C8（边缘混沌猜想）。
    \item \textbf{对C8的独立否证（v4新增）：}发现训练良好的Transformer的 $\lambda_{\max}$ 系统性显著偏离0——将否证C8的"边缘混沌"定位，尽管不一定否证不对称闭合的范畴分析本身。
\end{itemize}

上述经验证据——学术文献的独立验证、框架差异化预测的检验路径、工程案例的演示性支持、以及可操作的否证条件——联合构成了本文框架当前的经验支撑基础。同时，这些证据也暴露了多处尚未闭合的问题。以下逐项展开。


% ===================================================================
\section{未竟讨论}
\label{sec:open}

\subsection{数学基础的未闭合项}

\subsubsection{不对称闭合——从"弱闭合"到精确结构刻画}
\label{sec:open_asymmetric}

\textbf{本项在全框架中的地位变化：}在本文的早期版本中，残差加法与态射乘法之间的"不匹配"被描述为全框架最薄弱的环节——一个需要诚实标注的缺陷。经过第\ref{sec:residual_decomp}节的分析，这一判断已被根本性地修正。当前的理解如下。

\textbf{精确定位：}所谓"不匹配"是残差流张量积 $\mathbb{R}^n \otimes \mathbb{R}^d$ 中两个因子的范畴状态不对称：

\begin{table}[H]
\centering
\caption{不对称闭合的精确结构}
\begin{tabular}{@{}p{2.5cm}p{3cm}p{3.5cm}p{2cm}@{}}
\toprule
\textbf{因子} & \textbf{范畴状态} & \textbf{累积方式} & \textbf{米田} \\
\midrule
$\mathbb{R}^d$（语义空间） & 闭合——W矩阵可复合 & 乘性累积（$W_l \cdots W_1$） & $\checkmark$ \\
$\mathbb{R}^n$（路由空间） & 不闭合——A矩阵跨层不可复合 & 无累积（每层重新参数化） & $\times$ \\
整体残差流 & 加法连接（离散Euler积分器）——唯一能维持边缘混沌的累积方式 & — & — \\
\bottomrule
\end{tabular}
\end{table}

\textbf{关键推进：}不对称闭合不是需要修复的缺陷——它是Transformer能够产生语义结构的数学前提。
\begin{itemize}[leftmargin=*,itemsep=0.2em]
    \item $\mathbb{R}^n$ 闭合 + $\mathbb{R}^d$ 闭合 $\to$ Markov平稳分布 $\to$ 同质化 $\to$ 语义死亡
    \item $\mathbb{R}^n$ 不闭合 + $\mathbb{R}^d$ 不闭合 $\to$ 无语义自足结构 $\to$ 深度坍缩为单层
    \item $\mathbb{R}^n$ 不闭合 + $\mathbb{R}^d$ 闭合（实际架构）$\to$ 语义空间有自足结构 + 路由在演化中的语义空间上动态决策 $\to$ 语义动力学
\end{itemize}

\textbf{因果耦合的精确陈述：}W的复合改变语义空间坐标 $\to$ A依赖语义空间 $\to$ A逐层重新参数化 $\to$ A跨层不复合。W的复合成功 = A的复合失败的原因。两个因子的闭合状态不对称，且因果互补。

\textbf{残差加法的功能角色（重新理解）：}加法不只是"与态射复合不匹配"——它是唯一能长期维持边缘混沌（第\ref{sec:edge_chaos}节）的累积方式。乘法残差将导致轨线指数衰减到0（深度坍缩）或指数发散（表示爆炸）。加法残差通过 Jacobian $= I + \partial F/\partial x$ 使特征值维持在 $\approx 1$ 附近 $\to$ 边缘混沌可持续 $\to$ 深度有持续信息增量。

\textbf{两项独立研究的经验支撑（已在第\ref{sec:causal_coupling}节详述）：}Anthropic的Transformer Circuits框架 \cite{elhage2022transformer_circuits} 的路径展开定理和Deep Delta Learning \cite{zhang2025deep_delta} 的门控残差实验，分别从"加性何以能够工作"和"加性确实有限"两个方向为不对称闭合提供了独立支撑。

\textbf{当前闭合状态：}
\begin{itemize}[leftmargin=*,itemsep=0.2em]
    \item 结构描述：\textbf{严格层}——$\mathbb{R}^n$ 和 $\mathbb{R}^d$ 的范畴状态判定是数学事实
    \item 功能论证（不对称闭合 = 功能性的）：\textbf{形式论证层}——从范畴状态 + 反事实推理推导
    \item 强范畴论闭合：\textbf{未完成}——引入Effectus理论 \cite{cho2015effectus} 或Para构造 \cite{fong2019para} 以同时建模加性与乘性结构在数学上具有可行性，但复杂度远超出本文当前的形式化水平，且即便完成也仅能描述多层结构类型，无法给出收敛保证或表达能力的定量预测
\end{itemize}

\textbf{与早期版本的姿态差异：}早期版本把"弱闭合"当作需要诚实标注的\textbf{缺陷}。当前版本将其重新理解为Transformer架构的\textbf{功能性数学结构}——精确刻画了哪些部分闭合、哪些不闭合、为什么不对称、以及为什么这个不对称是语义动力学的前提。论证姿态从"我们这里不够好"变为"我们找到了为什么它必须是这样的"。

\subsubsection{有限范畴下"截断投影"的自洽性问题}

本文在第\ref{sec:truncated_projection}节中使用的数学设定——FinStoch范畴、有限token词汇表——恰好构成一个有限维范畴。在有限范畴上，分布语义嵌入本身即为有限维映射；"截断投影"在这一设定下作为原则上不可消除的信息损失来源，其论证基础存在自洽性缺口。修复方向为将"语义空间"显式拆分为两个层次：底层为由有限token集生成的FinStoch范畴（有限维，不需截断），上层为所有可能的概念和语义关系构成的连续语义空间（无限维）。截断发生在从上层连续语义空间到底层有限参数流形的映射过程中——此过程中的信息损失是真实的，但需要定义上层范畴与底层范畴之间的映射关系，这一工作尚未完成。

\subsubsection{交叉熵不蕴含函子性}

训练目标函数 $\mathcal{L} = -\sum \log P(\text{next\_token}|\text{context})$ 促使模型匹配训练数据中的条件概率分布，并不促使模型满足态射复合的结构保持条件 $F(g \circ f) = F(g) \circ F(f)$。交叉熵最小化与函子性之间在数学上不存在蕴含关系。当前修复策略为：不声称梯度下降在逻辑上蕴含函子性的获得；改为"训练收敛赋予映射行为在训练分布内的稳定性，经验观察表明在该稳定性条件下模型习得的映射恰好呈现出近似函子的结构保持性质"。这一表述将函子性从逻辑后承重新定位为经验发现。

\subsubsection{低秩限制与思维链后训练效果之间的互斥}

第\ref{sec:truncated_projection}节的核心依据——参数空间有效秩仅数百维量级——指向一个容量层面的诊断：有限维参数流形不足以编码完整的逻辑结构。然而，思维链强化学习后训练（CoT RL）在经验上被观察到系统性提升了模型的逻辑推理能力。两个命题在逻辑上不可同时为真：(A)若数百维有效秩确实构成逻辑结构表征的容量瓶颈，则CoT RL——其操作仍然在同一个低维参数空间内——不可能"补全"逻辑能力，其观察到的改善需被解释为在已有表征空间内的路径优化而非容量扩展；(B)若CoT RL确实补全了逻辑能力，则有效秩的容量瓶颈诊断不成立——瓶颈的性质需被重新解释为预训练目标（而非参数容量）未引导参数去表征逻辑结构。\textbf{当前未在这两条互斥路径之间做出选择，此张力构成本框架内部最尖锐的未解决矛盾。}本节的其余部分诚实标注：若路径B成立（容量足够而预训练目标未引导），则第\ref{sec:truncated_projection}节截断投影作为原则性限制的论证强度将被显著削弱——进而第\ref{sec:formal}节的四项推导（均以P1和P2为前提，而P1和P2源自截断投影约束）的可靠性将受到传导性影响。在路径选定之前，第\ref{sec:formal}节四项推导的可靠性以路径A为默认前提。

\subsubsection{Markov核对语义关系捕获范围的边界}

Markov核天然适配于捕获概率依赖关系——共现、联想、上下文替换等分布性语义关系属于可直接建模的范畴。组合语义（如"红苹果"的语义由"红"与"苹果"各自的语义映射经由复合产生）和结构化关系（部分-整体、因果等需要多个映射之间的交换约束来编码的关系）可通过范畴中的图表结构加以论证，但论证的具体构造尚未完成。规范性关系——"概率为1"与"逻辑必然"——在Markov范畴内不可区分，前者为度量概念而后者为证明论概念。视角依赖关系——指称、意向性、反事实条件句——依赖于认知主体的视角和模型无法访问的因果信息，无法被主体中性的Markov核完全捕获。这两类关系构成Markov核语义建模能力的原理性边界。

\subsubsection{边缘混沌假说的经验验证缺口（v4新增）}
\label{sec:open_edge_chaos}

第\ref{sec:edge_chaos}节的边缘混沌假说目前面临以下未闭合的经验问题：
\begin{itemize}[leftmargin=*,itemsep=0.2em]
    \item $\lambda_{\max} \approx 0$ 的预测（T8）尚未在任何已发表工作中得到直接验证。混沌三条件（路径依赖 + 非线性折叠 + 局部扩张）的"同时成立"判断来自对Transformer计算图的逐项数学分析，不是数值测量。
    \item 风险：实际 $\lambda_{\max}$ 可能系统性偏离0，此时需要解释为何Transformer仍能工作——可能"边缘"条件比 $\lambda_{\max} \approx 0$ 更宽松（例如仅要求 $\lambda_{\max}$ 在一个较宽的临界区间内而非精确等于0）。
    \item 风险：残差强度调节（T9）在真实模型中可能因LayerNorm的补偿效应而无法产生预测的分叉行为——LayerNorm在每层后重新缩放表示，可能吸收 $\varepsilon$ 变化的影响。
    \item 当前状态：诚实标注为概念构造层。T7--T10的检验结果将决定该假说是升级为形式论证层还是被否证。
\end{itemize}

\subsection{内部逻辑张力}

\subsubsection{动机缺失的根因未明}

本文框架观察到LLM Agent缺乏自发的、由内部状态驱动的知识获取动机，但尚未从框架的数学基础中为该缺失提供结构性推导。参数冻结假设可解释"模型不能在部署后更新其参数化知识"，但不能解释"模型为什么没有表现出基于好奇、认知不适或新奇感的自发求知行为"。在当前的论证状态下，动机缺失被诚实标注为一项经验观察——它与其他五项缺失（语义空间不完备、模态不对齐、无渐进记忆、顺序非时间、路由不复合）在论证强度上存在性质差异：前五项各自具有从框架的数学结构中出发的推导链，而动机缺失目前仅具有观察地位。

\subsubsection{多模态判断的适用范围}

本文在早期工作中提出的"其他模态仅是形式对齐或语言模拟"这一表述需接受范围限定。当前最新多模态模型（GPT-4o, Gemini 2.5系列）的原生多模态能力来自于文本与视觉/音频信号的联合训练，其机制超出了简单的表示空间对齐。需在以下两个层面之间做出区分：(1)多模态模型在工程基准上是否表现出有效的跨模态能力——此属工程事实层面的问题，当前证据倾向于肯定回答；(2)多模态模型在非语言模态中是否获得了与人类同构的因果语义结构——此属哲学主张层面的问题，本文的否定立场需要更精确的经验操作化。两个层面的讨论不应混为一体。

\subsection{发散框架的精度缺口}

收敛任务部分拥有从前提系统出发的形式化推导链（P0--P6, B1--B3，四项形式推导，一项三段论）。发散任务部分的核心洞察——条件概率采样的机制还原、偏好信号的信息论约束——已达到概念层面的精度，但整体的形式化密度尚未追平收敛部分。两个具体的开放问题：(1)偏好信号的信息论量化刻画——是否存在一个类似于验证信号覆盖度条件 $\verifspace = \taskcspace$ 的充分性判据，用于判定发散任务中偏好信号的信息量是否足以将采样过程驱动至满意的目标区域？(2)发散过程中观察到的新语义或新关联的涌现——当前框架将发散过程刻画为已有语义空间的偏好引导重组，但本文框架的早期表述暗示发散过程可能具有超出重组的生成性。若模型仅执行自回归条件概率采样，"全新语义"的涌现机制在计算层面如何被解释是一个尚未闭合的问题。第\ref{sec:divergent_formal}节的v4补注提出了一个方向性提示——新语义的涌现可能对应混沌轨线进入 $\mathbb{R}^d$ 中训练分布未曾覆盖的新吸引域——但此连接尚未展开为系统论证。

\textbf{[留空——待补充实证：发散任务过程中观察到的新语义涌现现象的系统记录与分析]}

\subsection{混合任务的复杂性}

本文仅从概念层面识别了发散与收敛的三种混合模式——序列混合（发散$\to$收敛$\to$发散$\to\cdots$的交替）、层级混合（高层处于发散模式而低层处于收敛模式）、并行混合（不同子任务在同一时刻处于不同模式）——但尚未系统分析每种模式下具体的Harness设计约束和任务完成率的结构下界。混合模式之间切换的代价量化——在不同模式间转换所消耗的上下文容量与操作者注意力资源——是尚未完成的理论工作。

\textbf{[留空——待补充实证：发散向收敛过渡的工程过程分析（本文网文创作框架中从结构大纲的发散设计到可验证章节标准的结晶过程即构成此类过渡的一个具体实例）]}

\subsection{"中期记忆不存在"论断的操作化}

当前论证——参数化记忆与上下文记忆分属不同数学结构，二者之间在本体论层面上不存在连续的过渡结构——具有明确的哲学含义，但若停留在本体论层面则难以被经验证据所否证。建议将该论断操作化为以下工程命题："在参数部署后冻结的标准Transformer部署范式下，上下文窗口中的信息无法获得参数化记忆所具有的收敛约束赋予的结构稳定性——无论采用何种上下文管理策略。"此操作化版本具备可检验性：若某一"中期记忆"工程方案能够展示上下文信息的跨会话时间性动态（衰减曲线、巩固效应、自发联想重组），且这些动态不由外部系统的显式规则驱动而由模型内部的计算结构自发产生，则该操作化命题在工程意义上被否证。

\textbf{补注——Fine-tuning与RAG推论的落点说明：}两项推论——"Fine-tuning不能修复上下文驱动分解决策的不可靠性"以及"RAG没有创造新的记忆结构类别"——均从记忆的结构性二分（\S\ref{sec:memory_dichotomy}）中直接推出。Fine-tuning改变的是参数空间中存储的长期记忆内容，无法触及"上下文中的信息本身不经收敛"这一结构事实——其更深层的根因在 $\mathbb{R}^n$ 跨层不复合（第\ref{sec:routing_space}节），微调不改变注意力机制的跨层范畴状态。RAG系统检索到的文段在注入上下文窗口的瞬间即成为短期记忆——一段被追加至序列中的文本，其"新鲜度"由外部索引的时间戳而非模型内部的记忆动态决定。RAG优化的是"何种信息在何种时机进入窗口"的调度策略，而非在参数空间结构与序列顺序结构之间创造出具有独立数学性质的第三种存储形式。

\subsection{Agent Scaling Law的精确预测形式}

C6猜想当前处于方向性预测阶段。使之成为严格的可检验科学假说，需要以下三方面的精确化：(1)平台期出现时的多步任务复杂度拐点的量化定位——在何种任务复杂度量级上，Agent多步性能的Scaling曲线开始与单步性能曲线出现统计学上显著的分化；(2)平台期差距的量级估计——两条曲线在平台期区域的性能差距的预期量级与模型规模、验证信号覆盖度之间的函数关系；(3)多步任务性能上界的函数形式——即 $\verifspace/\taskcspace$ 与性能上限之间的映射关系。

\subsection{开放研究方向}

\begin{enumerate}[leftmargin=*,itemsep=0.3em]
    \item \textbf{Agent可靠性理论的形式化：}将测试收敛、层级错误传播和根需求外在性整合为一个可操作的Agent可靠性度量框架。这是该框架中与当前工程实践对接最为直接、做出可检验预测最为便利的部分。
    \item \textbf{范畴论内核的形式化推进：}将分布语义嵌入与截断投影的关系从概念构造推进至具有可操作定义的数学对象。需与Categorical Invariants of Learning Dynamics \cite{categorical_invariants_2025}——当前唯一将学习定义在范畴论函子框架内的独立工作——进行实质性学术对话。
    \item \textbf{偏好信号的量化理论：}为发散任务建立与收敛任务中验证信号覆盖度条件 $\verifspace = \taskcspace$ 对等的分析工具。
    \item \textbf{架构改进的理论边界分析：}系统分析何种架构层面的改变可以移动本文识别的各项结构性边界——推理时参数持续学习、非softmax信息传递机制、具时间性衰减与巩固的记忆结构、可复合路由机制。
    \item \textbf{与记忆分类学的学术对话：}本文的"无中间过渡记忆"论断与2025年主流记忆研究（Hippocampus Revolution, Memory Taxonomy \cite{memory_taxonomy_2025}）构成了直接的理论冲突。需要将"不存在"的操作化定义明确写入学术对话框架，使冲突可被经验证据所裁决。
    \item \textbf{与多Agent理论的对接：}Stanford MACI/SocraSynth \cite{stanford_maci} 等九柱多Agent理论在Agent层级的覆盖面上具有广度优势，本文可在数学深度上提供互补。
    \item \textbf{边缘混沌假说的经验检验（v4新增）：}T7--T10的检验将决定该假说的命运。其中T10（强制复合$\to$同质化）是最强的否证实验——如果强制复合后模型表现不退化，则不对称闭合的核心命题被否证。
    \item \textbf{$\mathbb{R}^n$ 端路由的量化理论（v4新增）：}当前分析确立了 $\mathbb{R}^n$ 不复合是数学事实，但尚未给出路由偏离复合程度的量化度量。第\ref{sec:causal_coupling}节中提到的"A$^{(l+1)}$ 偏离CK条件的程度"是可能的理论方向——其量化将直接关联到语义聚类的强度和边缘混沌的Lyapunov谱。
\end{enumerate}


% ===================================================================
\section{结论}

本文从Transformer架构的数学结构出发，构建了一个关于LLM Agent能力边界的统一理论框架。全文的论证可浓缩为以下三条核心结论。

\textbf{第一，Agent的若干可观察失败模式是结构必然，而非工程不足。}上下文驱动分解决策的不可靠性、跨会话记忆的不连续性、在验证信号缺失条件下闭环的不可闭合性——这些失败模式的严重程度随模型能力（规模、数据、窗口长度）的提升而减轻，但其存在性根源于三个结构约束的共同作用：参数空间的有效维度限制（截断投影）、上下文信息不经训练收敛的梯度锻造、以及参数化记忆与序列化记忆之间数学结构的根本异质性。v4新增了更深层的第四项结构约束——残差流 $\mathbb{R}^n \otimes \mathbb{R}^d$ 中路由空间（$\mathbb{R}^n$）的跨层不复合，这是前三个约束中"上下文不经收敛"的数学机制根源。这四个约束不由模型规模的扩大而改变其性质。补偿路径在于Harness层面的外部机制设计——人工确认节点、外部持久化记忆系统、偏好信号采集界面——而非等待模型内部能力的量变积累。

\textbf{第二，Agent行为必须按评价模式区分为收敛与发散两个基本类别。}两种模式的驱动机制在信息论层面上根本不同：收敛任务由可自动执行的验证信号驱动（测试收敛），发散任务由操作者的比较性偏好信号驱动（偏好驱动迭代）。将两种模式混为一谈——以收敛执行的管理框架去处理发散探索，或以发散模式的开放性去替代收敛执行的验证标准——是Agent系统工程失效的首要结构性原因。

\textbf{第三，操作者在人机协作中的角色具有结构性而非过渡性。}该角色——根需求锚定、关键分叉仲裁、偏好信号注入、过渡时机确认——来自Agent架构的结构性缺口。这些缺口的存在性决定于当前Transformer范式的底层数学结构（参数冻结、softmax注意力、模态信号的形式对齐而非因果闭环、$\mathbb{R}^n$ 跨层不复合）。只要这些结构特征不发生根本性的改变，操作者的上述角色就不会随模型能力的量变而趋于消失。人机协作因而不是UX层面的过渡方案——它是Agent架构在当前数学结构下的构成性要素。

回到本文开篇的十七条使用原则。它们的有效性来自同一个根源：每一条原则都在一个特定的操作节点上，将"未经收敛约束的信息"与"需要收敛约束的决策"进行了正确的分离或匹配。认识到这些原则背后统一的数学结构，不仅解释了它们为何在实践中反复验证有效，也划定了它们有效性的适用边界——以及超出这些边界后，操作者的结构性判断何以不可替代。

\noindent\textbf{致谢：}感谢审稿Agent在本文多次迭代中提供的系统性、多维度的结构性反馈。

% ===================================================================
% 附录
% ===================================================================
\appendix

\section{推理链全图}

\begin{figure}[H]
\centering
\begin{tikzpicture}[
    node distance=0.6cm,
    box/.style={rectangle,draw,rounded corners,text width=5.0cm,align=center,minimum height=0.7cm,font=\footnotesize},
    arrow/.style={-{Stealth},thick}
]
\node[box,fill=lightgray] (decomp) {\textbf{残差流 $\mathbb{R}^n \otimes \mathbb{R}^d$}\\$\mathbb{R}^d$ 语义空间（范畴闭合）\\$\mathbb{R}^n$ 路由空间（跨层不复合）};
\node[box,below=of decomp,fill=lightred] (asym) {\textbf{不对称闭合}\\W复合成功 $\to$ A复合失败\\因果耦合};
\node[box,below=of asym] (dim) {参数有限维 $\to$ \textbf{截断投影}};
\node[box,below=of dim,fill=lightred] (conv) {\textbf{训练收敛} vs \textbf{上下文不经收敛}};
\node[box,below=of conv,fill=lightred] (mem) {\textbf{记忆二分：}空间结构 vs 顺序结构};
\node[box,below=of mem] (edge) {\textbf{边缘混沌}\\不复合$\to$路径依赖+非线性+扩张\\$\lambda_{\max} \approx 0$};
\node[box,below=of edge] (split) {$\verifspace = \emptyset$?（判据）};

\node[box,fill=lightblue,below left=0.8cm and 2.5cm of split] (yes) {\textbf{收敛任务}\\\textbf{测试收敛}};
\node[box,fill=lightblue,below right=0.8cm and 2.5cm of split] (no) {\textbf{发散任务}\\\textbf{偏好驱动迭代}};

\node[box,fill=lightblue,below=of yes] (decomp2) {分解不可靠（三段论）};
\node[box,fill=lightblue,below=of decomp2] (hier) {层级收敛 P0--P4};
\node[box,fill=lightblue,below=of hier] (root) {根需求外在（形式论证）};

\node[box,fill=lightblue,below=of no] (pref) {偏好函数唯一导向};
\node[box,fill=lightblue,below=of pref] (trans) {发散$\to$收敛过渡};

\node[box,fill=lightgray,below=2.5cm of root,xshift=2.5cm] (human) {\textbf{人的结构性角色}};
\node[box,fill=lightgray,below=of human] (eng) {\textbf{使用原则 + 可行域边界}};

\draw[arrow] (decomp) -- (asym);
\draw[arrow] (asym) -- (dim);
\draw[arrow] (dim) -- (conv);
\draw[arrow] (conv) -- (mem);
\draw[arrow] (mem) -- (edge);
\draw[arrow] (edge) -- (split);
\draw[arrow] (split) -| (yes);
\draw[arrow] (split) -| (no);
\draw[arrow] (yes) -- (decomp2);
\draw[arrow] (decomp2) -- (hier);
\draw[arrow] (hier) -- (root);
\draw[arrow] (no) -- (pref);
\draw[arrow] (pref) -- (trans);
\draw[arrow] (root.east) -- ++(1.2,0) |- (human);
\draw[arrow] (trans.east) -- ++(1.2,0) |- (human);
\draw[arrow] (human) -- (eng);
\end{tikzpicture}
\caption{全文推理链全图（v4）。灰框表示数学事实或综合结论，红框表示核心理论构建，蓝框表示概念构造。各节点间的箭头表示逻辑依赖关系。v4新增顶层节点（残差流分解$\to$不对称闭合）和边缘混沌节点。}
\end{figure}

\section{论证强度分层标准}

\begin{table}[H]
\centering
\begin{tabular}{@{}p{2cm}p{5.5cm}p{5cm}@{}}
\toprule
\textbf{层级} & \textbf{判定标准} & \textbf{本文对应内容} \\
\midrule
严格层 & 数学事实，无近似，"="而非"$\approx$" & \S\ref{sec:semantic_space}（W矩阵范畴闭合、单层注意力=Markov核、米田成立）；\S\ref{sec:routing_space}（跨层复合失败、米田不适用）；\S\ref{sec:causal_coupling}（因果耦合）；\S\ref{sec:convergence_def}（收敛操作化判据） \\
定义层 & 概念的形式规定，不涉及真值判断 & \S\ref{sec:definitions}.1（核心概念定义表） \\
形式论证层 & 从明示前提经有效推理步骤导出结论，前提与推理规则均已公开 & \S\ref{sec:causal_coupling}（不对称闭合作为功能条件的推导）；\S\ref{sec:truncated_projection}（截断投影$\to$覆盖边界$\to$分布外暴露）；\S\ref{sec:memory_dichotomy}（记忆二分）；\S\ref{sec:formal}.3--8（四项推导+三段论+人的角色推导）；\S\ref{sec:open_asymmetric}（不对称闭合的功能论证） \\
概念构造层 & 核心机制已达概念精度，但整体推导密度尚未达到形式论证层 & \S\ref{sec:edge_chaos}（边缘混沌——三条件分析为形式论证，Lyapunov谱分析为概念构造）；\S\ref{sec:divergent_formal}（发散框架）；\S\ref{sec:agi}（AGI汇聚论证）；\S\ref{sec:open_edge_chaos}（边缘混沌验证缺口） \\
类比层 & 概念对应关系成立但存在已知的结构不匹配，已在正文中诚实标注 & 第\ref{sec:residual_decomp}节的不对称闭合分析已将该层的残余范围大幅缩减——原"三层贝叶斯类比"中第二层的不匹配现已由不对称闭合提供精确结构刻画 \\
经验推论层 & 从理论框架中演绎出的可观察预测，或援引自独立来源的经验证据 & \S\ref{sec:empirical}（验证与预测） \\
\bottomrule
\end{tabular}
\end{table}

\textbf{v4论证强度升级要点：}
\begin{itemize}[leftmargin=*,itemsep=0.2em]
    \item \S\ref{sec:routing_space}（跨层不复合）从早期版本隐含在闭合分析中升级为独立且标注强度的子节
    \item \S\ref{sec:causal_coupling}（因果耦合与不对称闭合）为v4新增的严格层+形式论证层内容
    \item \S\ref{sec:edge_chaos}（边缘混沌）为v4新增的概念构造层内容
    \item \S\ref{sec:open_asymmetric}（不对称闭合）从早期版本的"诚实标注缺陷"升级为"精确结构刻画+功能论证"
    \item T7--T10为v4新增的经验推论层内容
\end{itemize}

\section{框架推论九条在本文中的具体落点}

\begin{table}[H]
\centering
\begin{tabular}{@{}c@{}p{8.5cm}@{}}
\toprule
\textbf{推论} & \textbf{本文位置} \\
\midrule
1. 幻觉是结构必然非工程bug & \S\ref{sec:truncated_projection}（截断投影与收敛覆盖边界）；\S\ref{sec:definitions}.1（分布外暴露风险的形式定义） \\
2. Fine-tuning不修复上下文不可靠 & \S\ref{sec:open}.5补注（记忆二分$\to$微调不触及上下文不经收敛的结构事实——更深根因在 $\mathbb{R}^n$ 跨层不复合） \\
3. RAG不创造新记忆结构 & \S\ref{sec:memory_dichotomy}（RAG管理上下文容量）；\S\ref{sec:open}.5补注 \\
4. 可验证性密度的任务域阶梯 & \S\ref{sec:empirical}.1表格（代码$>$数学$>$翻译$>$开放写作） \\
5. 多Agent辩论的验证信号机制 & \S\ref{sec:empirical}.1表格与相关文献引用 \\
6. 过度思考的逆缩放预期 & \S\ref{sec:empirical}.1；T5（\S\ref{sec:predictions}） \\
7. 人介入的非过渡性 & \S\ref{sec:formal}.8（人的结构性角色推导） \\
8. 自主性-可靠性的结构权衡 & C5（\S\ref{sec:definitions}.2）；T6（\S\ref{sec:predictions}） \\
9. Agent任务Scaling的先行平台期 & C6（\S\ref{sec:definitions}.2）；T1（\S\ref{sec:predictions}）；\S\ref{sec:open}.6 \\
\bottomrule
\end{tabular}
\end{table}

\section{草稿57条与本文各节对应关系}

\begin{table}[H]
\centering
\begin{tabular}{@{}c@{}p{4.5cm}p{5cm}@{}}
\toprule
\textbf{草稿条目} & \textbf{本文对应章节} & \textbf{性质} \\
\midrule
1--6（数学基础） & \S\ref{sec:semantic_space}, \S\ref{sec:routing_space}, \S\ref{sec:causal_coupling}, \S\ref{sec:truncated_projection} & 严格层+形式论证层 \\
7--9（逻辑模拟+泄露+CoT RL） & \S\ref{sec:truncated_projection}, \S\ref{sec:open}.1.4, \S\ref{sec:open} & 分散至公理+未闭合 \\
10--11（基本原理与chat原则） & \S\ref{sec:convergence_def}, \S1.1（P1--P2） & 桥梁性断言 \\
12--17（Agent特征与记忆） & \S\ref{sec:memory_dichotomy}, \S\ref{sec:formal}.1（P0+P1） & 形式论证层 \\
18--23（收敛任务） & \S\ref{sec:formal}.1--7 & 形式论证层（全文最成熟部分） \\
24--28（发散任务） & \S\ref{sec:divergent_formal}, \S\ref{sec:agi} & 概念构造层 \\
29--33（人的结构性角色） & \S\ref{sec:formal}.8 & 形式论证层 \\
34（真实任务混合性质） & \S\ref{sec:open}.4 & 经验推论层 \\
35（子Agent扁平极限） & \S1.1（P13）, \S\ref{sec:formal}.7 & 经验推论+形式论证 \\
36（社会类比） & \S\ref{sec:formal}.8（以结构推导替代社会类比） & 辅助框架理解 \\
37（操作者上限） & \S1.1（P14）, \S\ref{sec:formal}.8 & 原则+论证 \\
38（合作者能力匹配） & \S\ref{sec:formal}.8 & 论证层 \\
39（原则引出） & \S1.1--1.2 & 原则引出 \\
40--56（十七条使用原则） & \S1.1 & 经验归纳 \\
57（AGI不可达） & C7（\S\ref{sec:definitions}.2）, \S\ref{sec:agi} & 猜想+汇聚论证 \\
框架推论九条 & 见附录推论落点表 & 分散至各论证层 \\
硬伤分析+闭合分析 & \S\ref{sec:open}.1 & 未竟讨论 \\
人机协作洞察 & \S\ref{sec:divergent_formal}（直觉/逻辑/训练-推理边界）, \S\ref{sec:formal}.8（收敛放手/发散不可脱手） & 分散至发散+角色 \\
\bfseries v4新增：笔记01--10 & \S\ref{sec:residual_decomp}, \S\ref{sec:edge_chaos}, \S\ref{sec:predictions}（T7--T10）, \S\ref{sec:open_asymmetric}, \S\ref{sec:open_edge_chaos} & 分散至严格层+概念构造层+经验推论层+未闭合 \\
\bottomrule
\end{tabular}
\end{table}

% ===================================================================
\newpage
\printbibliography

\end{document}
